\chapter[Canonical Geometry and Reflexive Polytopes]
{Canonical Geometry and\\ Reflexive Polytopes}
\label{chap:canonical-geometry-reflexive-polytopes}

\begin{chapteroverview}
Chapter~4 gave the striking formula
\[
K_{X_\Sigma}=-\sum_{\rho\in\Sigma(1)}D_\rho.
\]
At the time, we used it mainly as a divisor calculation. We now ask
what this formula means geometrically and what it can do.

The starting point is the invariant differential
\[
\frac{d\character{m}}{\character{m}}.
\]
It is regular on the dense torus, but it acquires logarithmic poles
when the torus is compactified by its boundary divisors. The residues
of those poles are exactly the pairings
\[
\langle m,u_\rho\rangle.
\]
Thus the character lattice, the toric boundary, and differential forms
fit into one dictionary.

The rest of the chapter will turn to the anticanonical divisor
\[
-K_{X_\Sigma}=\sum_\rho D_\rho,
\]
its positivity, and the lattice polytope
\[
P_{-K_X}
=
\{m\in M_\RR:\langle m,u_\rho\rangle\geq-1\text{ for every }\rho\}.
\]
This will lead to Fano toric varieties, polar duality, and reflexive
polytopes.
\end{chapteroverview}

\begin{threedictionaries}
The canonical geometry of a toric variety begins with one simple
operation: logarithmically differentiating a character.

\medskip

\begin{center}
\small
\begin{tabularx}{\linewidth}{@{}
  >{\raggedright\arraybackslash}X
  >{\raggedright\arraybackslash}X
  >{\raggedright\arraybackslash}X@{}}
\toprule
\textbf{Algebra}
&
\textbf{Geometry}
&
\textbf{Combinatorics}
\\
\midrule
character \(\character{m}\)
&
invariant logarithmic form
\(d\character{m}/\character{m}\)
&
lattice vector \(m\in M\)
\\
pole along \(D_\rho\)
&
residue on a boundary component
&
pairing \(\langle m,u_\rho\rangle\)
\\
top exterior power
&
canonical sheaf
&
negative of the toric boundary
\\
anticanonical sections
&
Fano geometry
&
lattice points and polar polytopes
\\
\bottomrule
\end{tabularx}
\end{center}
\end{threedictionaries}

\begin{surfacelaboratory}
For a smooth complete toric surface, the boundary is a cycle of smooth
rational curves. The identity
\[
K_X+\partial X=0
\]
says that the pair \((X,\partial X)\) carries a nowhere-vanishing
logarithmic two-form.

Later, the polygon of \(-K_X\) will let us see whether the surface is
Fano, compute \((-K_X)^2\), and recognize how toric blowups change the
anticanonical geometry.
\end{surfacelaboratory}

\section{Logarithmic forms and the toric boundary}
\label{sec:logarithmic-forms-toric-boundary}
\index{logarithmic differential form}
\index{toric boundary}
\index{logarithmic cotangent bundle}
\index{residue sequence}

In \cref{sec:canonical-divisor-singularity-criteria}, we computed the
canonical divisor by examining one invariant top form. We now widen
the lens. Instead of studying only the top exterior power, we will
understand all invariant one-forms and how they interact with every
boundary component.

Throughout the main discussion, \(X_\Sigma\) is smooth of dimension
\(n\), and
\[
\boundaryX
=
\sum_{\rho\in\Sigma(1)}D_\rho
\]
is its reduced toric boundary. The singular case will be recorded at
the end as an optional extension.

\begin{agcheatsheet}
Let \(X\) be smooth and let
\[
D=D_1+\cdots+D_r
\]
be a divisor with simple normal crossings. Near a point of \(X\), one
can choose local coordinates
\[
x_1,\dots,x_n
\]
so that
\[
D=V(x_1\cdots x_s)
\]
for some \(s\leq n\).

Ordinary regular one-forms are locally generated by
\[
dx_1,\dots,dx_n.
\]
A \emph{logarithmic one-form along \(D\)} is allowed to have a simple
pole in each boundary direction. Such forms are locally generated by
\[
\frac{dx_1}{x_1},\dots,\frac{dx_s}{x_s},
\quad
dx_{s+1},\dots,dx_n.
\]
They form a locally free sheaf denoted
\[
\Omega_X^1(\log D).
\]

The coefficient of \(dx_i/x_i\), restricted to \(D_i\), is the
\emph{residue} along \(D_i\). Thus
\[
\operatorname{Res}_{D_i}\!\left(\frac{dx_i}{x_i}\right)=1.
\]
A logarithmic pole is therefore controlled: it remembers one regular
function on the divisor where the pole occurs.
\end{agcheatsheet}

\subsection{Logarithmically differentiating a character}

Let \(m\in M\). On the dense torus \(T_N\), the character
\(\character{m}\) is an invertible regular function, so the expression
\[
\operatorname{dlog}(\character{m})
:=
\frac{d\character{m}}{\character{m}}
\]
is a regular one-form on \(T_N\).

The notation is justified by the familiar rule
\[
\operatorname{dlog}(fg)
=
\operatorname{dlog}(f)+\operatorname{dlog}(g).
\]
Since
\[
\character{m+m'}=\character{m}\character{m'},
\]
we obtain
\[
\operatorname{dlog}(\character{m+m'})
=
\operatorname{dlog}(\character{m})
+
\operatorname{dlog}(\character{m'}).
\]
Thus logarithmic differentiation turns multiplication of characters
into addition of differential forms.

Choose a basis \(m_1,\dots,m_n\) of \(M\), and write
\[
t_i=\character{m_i}.
\]
If
\[
m=a_1m_1+\cdots+a_nm_n,
\]
then
\[
\character{m}=t_1^{a_1}\cdots t_n^{a_n},
\]
and hence
\[
\boxed{
\operatorname{dlog}(\character{m})
=
\sum_{i=1}^n a_i\frac{dt_i}{t_i}.
}
\]

\begin{proposition}[Invariant one-forms on the torus]
\label{prop:invariant-one-forms-torus}
The assignment
\[
m\longmapsto
\operatorname{dlog}(\character{m})
\]
induces an isomorphism
\[
M\otimes_\ZZ\Bbbk
\simeq
H^0(T_N,\Omega^1_{T_N})^{T_N}
\]
from the character lattice tensored with \(\Bbbk\) to the vector
space of torus-invariant regular one-forms.
\end{proposition}

\begin{proof}
For a basis \(m_1,\dots,m_n\) of \(M\), the torus is
\[
T_N\simeq\Gm^n
\]
with coordinates \(t_1,\dots,t_n\). Its regular one-forms form the
free \(\Bbbk[M]\)-module generated by
\[
\frac{dt_1}{t_1},\dots,\frac{dt_n}{t_n}.
\]
These generators are invariant under translation by the torus. An
invariant form cannot have a nonconstant Laurent polynomial as a
coefficient, so the invariant forms are precisely their constant
linear combinations. The displayed formula identifies those
combinations with \(M\otimes_\ZZ\Bbbk\).
\end{proof}

\noindent
\textit{Reference.}
This is the torus calculation at the beginning of
Cox--Little--Schenck, \S8.1, especially equation~(8.1.1).

\begin{observation}
The representation-theoretic language is optional here. The torus
acts on itself by translation, hence on differential forms by
pullback. The forms
\[
\frac{dt_i}{t_i}
\]
are fixed by this action. Thus the lattice \(M\) gives a preferred
basis of invariant cotangent directions everywhere on the torus.
\end{observation}

\subsection{What happens at the boundary}

A logarithmic form is designed to survive when the torus is completed
by boundary divisors. The local toric chart makes this completely
visible.

Let \(\sigma\in\Sigma\) be a smooth cone of dimension \(r\). Choose a
basis
\[
u_1,\dots,u_n
\]
of \(N\) such that
\[
\sigma=\Cone{u_1,\dots,u_r}.
\]
Let \(m_1,\dots,m_n\) be the dual basis, and set
\[
x_i=\character{m_i}.
\]
Then
\[
U_\sigma
\simeq
\AAA^r\times\Gm^{n-r},
\]
and the toric boundary on this chart is
\[
\boundaryX\cap U_\sigma
=
V(x_1\cdots x_r).
\]
The divisor associated to the ray \(\rho_i=\RR_{\geq0}u_i\) is
\[
D_{\rho_i}\cap U_\sigma=V(x_i).
\]

\begin{figure}[ht]
\centering
\begin{tikzpicture}[scale=0.95]
  % Fan panel
  \begin{scope}[xshift=-4.15cm]
    \fill[DictionaryBlue!10] (0,0)--(2.45,0)--(0,2.45)--cycle;
    \draw[->,black!65] (-0.3,0)--(2.85,0) node[below right] {\(N_\RR\)};
    \draw[->,black!65] (0,-0.3)--(0,2.85);
    \draw[ray,DictionaryBlue] (0,0)--(2.35,0)
      node[below right] {\(u_1\)};
    \draw[ray,SurfaceGreen] (0,0)--(0,2.35)
      node[above left] {\(u_2\)};
    \node[font=\bfseries] at (1.2,3.08) {smooth cone \(\sigma\)};
    \node[DictionaryBlue] at (1.55,0.38) {\(\rho_1\)};
    \node[SurfaceGreen] at (0.38,1.55) {\(\rho_2\)};
  \end{scope}

  \draw[->,very thick,black!45] (-1.15,1.2)--(1.15,1.2)
    node[midway,above,black!60] {affine chart};

  % Affine chart panel
  \begin{scope}[xshift=4.15cm]
    \node[font=\bfseries] at (1.25,2.95) {\(U_\sigma\simeq\AAA^2\)};
    \fill[LightGray] (0,0) rectangle (2.5,2.3);
    \draw[->,black!65] (-0.18,0)--(2.9,0) node[below right] {\(x_1\)};
    \draw[->,black!65] (0,-0.18)--(0,2.65) node[above left] {\(x_2\)};
    \draw[DictionaryBlue,line width=2.4pt] (0,0)--(0,2.3);
    \draw[SurfaceGreen,line width=2.4pt] (0,0)--(2.5,0);
    \node[DictionaryBlue,rotate=90,fill=white,inner sep=1pt]
      at (-0.31,1.15) {\(D_{\rho_1}\)};
    \node[SurfaceGreen,fill=white,inner sep=1pt]
      at (1.25,-0.31) {\(D_{\rho_2}\)};
    \node[align=center,black!65,fill=white,rounded corners,inner sep=3pt]
      at (1.32,1.2) {dense torus \(T_N\)\\\(x_1x_2\neq0\)};
  \end{scope}
\end{tikzpicture}
\caption{On a smooth toric chart, rays become coordinate boundary
components. The form \(\frac{dx_1}{x_1}\wedge\frac{dx_2}{x_2}\) has a simple pole along
each colored boundary axis.}
\label{fig:smooth-cone-logarithmic-chart}
\end{figure}

For a general lattice vector
\[
m=a_1m_1+\cdots+a_nm_n,
\]
we have
\[
\operatorname{dlog}(\character{m})
=
\sum_{i=1}^n a_i\frac{dx_i}{x_i}.
\]
Along \(D_{\rho_i}=V(x_i)\), the residue is therefore \(a_i\). But
\[
a_i=\langle m,u_i\rangle.
\]
This gives the local formula
\[
\boxed{
\operatorname{Res}_{D_\rho}
\left(\frac{d\character{m}}{\character{m}}\right)
=
\langle m,u_\rho\rangle.
}
\]

\begin{pauseandthink}
Why is the primitive generator \(u_\rho\), rather than an arbitrary
nonzero lattice point on the ray, the correct vector in the residue
formula?

Compare this with the order formula
\[
\ord_{D_\rho}(\character{m})
=
\langle m,u_\rho\rangle.
\]
What is the common local calculation behind the two statements?
\end{pauseandthink}

\subsection{The logarithmic cotangent bundle is trivial}

The local forms
\[
\frac{dx_1}{x_1},\dots,\frac{dx_n}{x_n}
\]
are not an accidental coordinate basis. They are the forms attached
to the dual lattice basis \(m_1,\dots,m_n\). Changing the lattice
basis changes these forms by the same integral matrix. Consequently,
they patch canonically over all toric charts.

\begin{theorem}[Logarithmic forms on a smooth toric variety]
\label{thm:logarithmic-cotangent-trivial-toric}
Let \(X_\Sigma\) be a smooth toric variety with toric boundary
\(\boundaryX\). Then
\[
\boxed{
\Omega^1_{X_\Sigma}(\log\boundaryX)
\simeq
M\otimes_\ZZ\OX.
}
\]
Under this isomorphism,
\[
m\otimes1
\longleftrightarrow
\frac{d\character{m}}{\character{m}}.
\]
In particular, the logarithmic cotangent bundle is trivial.
\end{theorem}

\noindent
\textit{Reference.}
This is Cox--Little--Schenck, equation~(8.1.5), together with the
local logarithmic-form construction preceding it.

\begin{proof}
On a smooth chart
\[
U_\sigma\simeq\AAA^r\times\Gm^{n-r},
\]
the boundary has equation \(x_1\cdots x_r=0\). Hence
\(\Omega^1(\log\boundaryX)\) is freely generated by
\[
\frac{dx_1}{x_1},\dots,\frac{dx_r}{x_r},
\frac{dx_{r+1}}{x_{r+1}},\dots,
\frac{dx_n}{x_n}.
\]
For \(i>r\), the coordinate \(x_i\) is a unit, so
\(dx_i/x_i\) is an ordinary regular form. Sending \(m_i\) to
\(dx_i/x_i\) gives an isomorphism
\[
M\otimes_\ZZ\mathcal O_{U_\sigma}
\simeq
\Omega^1_{U_\sigma}(\log\boundaryX).
\]
The construction is expressed entirely through characters, so the
local isomorphisms agree on overlaps and patch globally.
\end{proof}

The residues assemble into an exact sequence.

\begin{theorem}[The toric residue sequence]
\label{thm:toric-residue-sequence}
For a smooth toric variety,
\[
0
\longrightarrow
\Omega^1_{X_\Sigma}
\longrightarrow
M\otimes_\ZZ\OX
\stackrel{\operatorname{Res}}{\longrightarrow}
\bigoplus_{\rho\in\Sigma(1)}\mathcal O_{D_\rho}
\longrightarrow0
\]
is exact, where
\[
\operatorname{Res}(m\otimes f)
=
\bigl(
\langle m,u_\rho\rangle f|_{D_\rho}
\bigr)_\rho.
\]
\end{theorem}

\noindent
\textit{Reference.}
This is Cox--Little--Schenck, Theorem~8.1.1, interpreted through
logarithmic forms and the Poincar\'e residue sequence~(8.1.6).

\begin{proof}[Conceptual proof sketch]
On a smooth affine chart, the boundary is a union of coordinate
hyperplanes. A logarithmic form lies in the ordinary cotangent sheaf
exactly when all of its logarithmic residues vanish. Conversely, one
may prescribe an arbitrary regular residue along each coordinate
boundary component by choosing the corresponding coefficient of
\(dx_i/x_i\). This proves exactness locally, hence globally.
\end{proof}

The middle term is simple, but the three terms encode different
parts of the geometry:

\begin{center}
\small
\renewcommand{\arraystretch}{1.25}
\begin{tabularx}{0.96\linewidth}{@{}
  >{\centering\arraybackslash}X
  >{\centering\arraybackslash}X
  >{\centering\arraybackslash}X@{}}
ordinary differentials
&
lattice-valued logarithmic differentials
&
residues on the boundary
\\
\(\Omega_X^1\)
&
\(M\otimes\mathcal O_X\)
&
\(\displaystyle\bigoplus_\rho\mathcal O_{D_\rho}\)
\\[-1mm]
\multicolumn{3}{c}{\(
\Omega_X^1
\longrightarrow
M\otimes\mathcal O_X
\stackrel{\operatorname{Res}}{\longrightarrow}
\displaystyle\bigoplus_\rho\mathcal O_{D_\rho}
\)}
\end{tabularx}
\end{center}

\subsection{Taking the top exterior power}

Choose a basis \(m_1,\dots,m_n\) of \(M\). The wedge
\[
\Omega_{\log}
=
\frac{d\character{m_1}}{\character{m_1}}
\wedge\cdots\wedge
\frac{d\character{m_n}}{\character{m_n}}
\]
is a nowhere-vanishing global section of
\[
\Omega^n_{X_\Sigma}(\log\boundaryX).
\]
A different lattice basis changes this form only by determinant
\(\pm1\). Thus the logarithmic canonical bundle is canonically
trivial up to the choice of orientation of \(M\).

For a simple-normal-crossing divisor \(D\) on a smooth variety,
\[
\Omega_X^n(\log D)
\simeq
\omega_X(D).
\]
Applying this to the toric boundary gives
\[
\omega_{X_\Sigma}(\boundaryX)
\simeq
\mathcal O_{X_\Sigma}.
\]
Equivalently,
\[
\boxed{
K_{X_\Sigma}+\boundaryX\sim0.
}
\]
In fact, using the particular rational form \(\Omega_{\log}\), we
recover the representative already proved in
\cref{thm:canonical-divisor-toric}:
\[
\operatorname{div}(\Omega_{\log})
=
-\sum_{\rho\in\Sigma(1)}D_\rho.
\]

\begin{observation}[The toric pair is logarithmically Calabi--Yau]
A smooth pair \((X,D)\) is often called \emph{log Calabi--Yau} when
\[
K_X+D\sim0.
\]
Thus every smooth toric variety, together with its full toric
boundary, is log Calabi--Yau.

This does \emph{not} mean that \(K_X\sim0\). The boundary is essential:
it is exactly what compensates for the poles of the invariant volume
form.
\end{observation}

\subsection{Two compactifications of a torus}

\begin{example}[The two ends of \(\PP^1\)]
The dense torus of \(\PP^1\) has coordinate \(t\). The invariant form
\[
\frac{dt}{t}
\]
has a simple pole at \(t=0\). Near infinity, use \(s=1/t\). Then
\[
\frac{dt}{t}
=
-\frac{ds}{s},
\]
so there is also a simple pole at \(s=0\), which is the point at
infinity. Therefore
\[
\operatorname{div}\!\left(\frac{dt}{t}\right)
=
-[0]-[\infty].
\]
The two boundary points are the two ends added when \(\Gm\) is
compactified to \(\PP^1\).
\end{example}

\begin{example}[The three boundary lines of \(\PP^2\)]
Let
\[
D_i=V(X_i)\subseteq\PP^2.
\]
On the torus chart \(X_0X_1X_2\neq0\), set
\[
x=\frac{X_1}{X_0},
\qquad
y=\frac{X_2}{X_0}.
\]
The invariant volume form is
\[
\Omega_{\log}
=
\frac{dx}{x}\wedge\frac{dy}{y}.
\]
Its simple poles along \(D_1\) and \(D_2\) are visible from
\(x=0\) and \(y=0\).

To inspect \(D_0\), work on \(X_1\neq0\) with
\[
u=\frac{X_0}{X_1}=\frac1x,
\qquad
v=\frac{X_2}{X_1}=\frac yx.
\]
Then
\[
\frac{dx}{x}=-\frac{du}{u},
\qquad
\frac{dy}{y}=\frac{dv}{v}-\frac{du}{u},
\]
and hence
\[
\Omega_{\log}
=
-\frac{du}{u}\wedge\frac{dv}{v}.
\]
Thus \(u=0\), which is \(D_0\), also contributes a simple pole. We
recover
\[
K_{\PP^2}=-(D_0+D_1+D_2)\sim-3H.
\]
\end{example}

\begin{pauseandthink}
On \(\PP^1\times\PP^1\), choose torus coordinates \(s,t\) and analyze
\[
\frac{ds}{s}\wedge\frac{dt}{t}.
\]
Locate its four poles and recover
\[
K_{\PP^1\times\PP^1}
\sim
-2F_1-2F_2.
\]
\end{pauseandthink}

\begin{observation}[What changes in the singular case]
For a normal toric variety, the ordinary cotangent sheaf need not be
locally free. One replaces it by the reflexive sheaf of Zariski
one-forms. Cox--Little--Schenck, Theorem~8.1.4, gives the corresponding
residue sequence, and Theorem~8.2.3 gives
\[
\omega_{X_\Sigma}
\simeq
\mathcal O_{X_\Sigma}\!\left(-\sum_\rho D_\rho\right).
\]

No singular differential-form technology is needed for the next
sections. The codimension-one formula from Chapter~4 remains the
working tool.
\end{observation}

The boundary has now appeared in two complementary roles:
\[
\boxed{
\begin{aligned}
\text{as poles of invariant forms:}
&\qquad K_X=-\boundaryX,\\
\text{as a positive divisor:}
&\qquad -K_X=\boundaryX.
\end{aligned}
}
\]
The next section studies when this positive-looking divisor is
actually ample.

\begin{guidingquestion}
The anticanonical divisor of a toric variety is the sum of all boundary
components:
\[
-K_{X_\Sigma}=\sum_\rho D_\rho.
\]
When does this divisor define an embedding into projective space?
What should its support function and its polytope look like when that
happens?
\end{guidingquestion}
\section{The anticanonical divisor and Fano varieties}
\label{sec:anticanonical-fano}
\index{anticanonical divisor}
\index{anticanonical polytope}
\index{Fano toric variety}
\index{Fano criterion}

The previous section explained why the toric boundary is anticanonical:
\[
-K_{X_\Sigma}=\sum_{\rho\in\Sigma(1)}D_\rho.
\]
The coefficients are all positive, but that alone does not make the
divisor positive in the geometric sense. The boundary components may
interact in a way that makes their sum ample, merely basepoint free, or
not even nef.

This section turns that question into a picture in \(M_\RR\).

\begin{agcheatsheet}
Let \(X\) be a smooth projective variety. The variety is called
\emph{Fano} if its anticanonical divisor is ample:
\[
-K_X\text{ is ample}.
\]
Equivalently, some positive tensor power of
\[
\omega_X^{-1}=\mathcal O_X(-K_X)
\]
gives a closed embedding into projective space.

For a normal variety, one must first ask whether \(K_X\) is Cartier.
A normal variety is called \emph{Gorenstein Fano} if
\[
-K_X\text{ is Cartier and ample}.
\]
On a smooth variety every Weil divisor is Cartier, so the two
expressions describe the same condition.
\end{agcheatsheet}

\noindent
\textit{Reference.}
This is Cox--Little--Schenck, Definition~8.3.1.

\subsection{The anticanonical polytope}

Since every coefficient of \(-K_{X_\Sigma}\) is equal to \(1\), its
divisor polytope has the especially uniform form
\[
\boxed{
P_{-K_X}
=
\left\{
 m\in M_\RR:
 \langle m,u_\rho\rangle\geq-1
 \text{ for every }\rho\in\Sigma(1)
\right\}.
}
\]
Each ray contributes a half-space whose boundary is one lattice unit
from the origin.

The origin satisfies every inequality strictly:
\[
\langle 0,u_\rho\rangle=0>-1.
\]
Thus \(0\) lies in the interior of \(P_{-K_X}\). When \(\Sigma\) is
complete, the polytope is bounded by
\cref{prop:complete-fan-bounded-PD}.

\begin{proposition}[The distinguished interior lattice point]
\label{prop:anticanonical-unique-interior-point}
Let \(X_\Sigma\) be complete. Then the origin is the unique interior
lattice point of \(P_{-K_X}\):
\[
\operatorname{Int}(P_{-K_X})\cap M=\{0\}.
\]
\end{proposition}

\begin{proof}
Suppose \(m\in M\) lies in the interior. Then
\[
\langle m,u_\rho\rangle>-1
\]
for every ray. The pairing is integral, so
\[
\langle m,u_\rho\rangle\geq0
\qquad
\text{for every }\rho.
\]
Since the cones of a complete fan cover \(N_\RR\), the linear
functional \(m\) is nonnegative in every direction. Applying this to
both \(u\) and \(-u\) shows that \(m=0\).
\end{proof}

This conclusion does not yet say that \(-K_X\) is ample. The
inequalities may be redundant, the vertices may fail to be lattice
points, or the normal fan of the polytope may be coarser than the fan
we started with.

\subsection{The toric Fano criterion}

Assume now that \(X_\Sigma\) is smooth and complete. Then \(-K_X\) is
Cartier, and its support function satisfies
\[
\psi_{-K_X}(u_\rho)=-1
\qquad
\text{for every ray }\rho.
\]

\begin{theorem}[The toric Fano criterion]
\label{thm:toric-fano-criterion}
For a smooth complete toric variety \(X_\Sigma\), the following are
equivalent:

\begin{enumerate}[label=(\alph*),wide=0pt]
\item \(X_\Sigma\) is Fano.

\item The support function \(\psi_{-K_X}\) is strictly convex.

\item The anticanonical polytope \(P_{-K_X}\) is a lattice polytope
whose normal fan is exactly \(\Sigma\).

\item For every maximal cone \(\sigma\), its anticanonical Cartier
datum \(m_\sigma\) is a vertex of \(P_{-K_X}\), and the normal cone at
that vertex is \(\sigma\).
\end{enumerate}
\end{theorem}

\begin{proof}
The equivalence of \textup{(a)} and \textup{(b)} is
\cref{thm:ample-strictly-convex} applied to \(-K_X\).

When a Cartier divisor is ample, its support function is strictly
convex, its Cartier data are the vertices of its polytope, and the
polytope has normal fan \(\Sigma\). Conversely, if the normal fan of
\(P_{-K_X}\) is \(\Sigma\), then the support function has a genuinely
different linear piece on every maximal cone, hence is strictly
convex. This gives the remaining equivalences.
\end{proof}

\begin{corollary}[Facets at lattice distance one]
\label{cor:fano-facets-distance-one}
If \(X_\Sigma\) is smooth and Fano, then every ray inequality
\[
\langle m,u_\rho\rangle\geq-1
\]
cuts out a facet of \(P_{-K_X}\). Its primitive inward normal is
\(u_\rho\), so every facet is exactly one lattice unit from the
origin.
\end{corollary}

\begin{threedictionaries}
For a smooth complete toric variety, Fano geometry is the same
condition viewed in three languages.

\medskip

\begin{center}
\small
\begin{tabularx}{\linewidth}{@{}
  >{\raggedright\arraybackslash}X
  >{\raggedright\arraybackslash}X
  >{\raggedright\arraybackslash}X@{}}
\toprule
\textbf{Algebra}
&
\textbf{Geometry}
&
\textbf{Combinatorics}
\\
\midrule
anticanonical ring
\(\bigoplus_{k\geq0}H^0(X,-kK_X)\)
&
anticanonical embeddings for \(k\gg0\)
&
lattice points of \(kP_{-K_X}\)
\\
\(-K_X=\sum_\rho D_\rho\)
&
\(X\) is Fano
&
\(\psi_{-K_X}\) is strictly convex
\\
Cartier data \(m_\sigma\)
&
distinct local anticanonical frames
&
vertices of a polytope with normal fan \(\Sigma\)
\\
\bottomrule
\end{tabularx}
\end{center}
\end{threedictionaries}

\subsection{The first anticanonical polygons}

\begin{example}[The anticanonical triangle of \(\PP^2\)]
For the rays
\[
u_0=-e_1-e_2,
\qquad
u_1=e_1,
\qquad
u_2=e_2,
\]
the anticanonical polygon is
\[
P_{-K_{\PP^2}}
=
\left\{
(a,b):a\geq-1,\ b\geq-1,\ a+b\leq1
\right\}.
\]
Hence
\[
P_{-K_{\PP^2}}
=
\Conv\{(-1,-1),(2,-1),(-1,2)\}.
\]
All three inequalities are active, and its normal fan is the fan of
\(\PP^2\). Thus \(-K_{\PP^2}\sim3H\) is ample.
\end{example}

\begin{example}[The anticanonical square of
\(\PP^1\times\PP^1\)]
The four ray generators are \(\pm e_1,\pm e_2\), so
\[
P_{-K_{\PP^1\times\PP^1}}
=
[-1,1]\times[-1,1].
\]
Its normal fan is the product fan, and
\[
-K_{\PP^1\times\PP^1}
\sim2F_1+2F_2
\]
is ample.
\end{example}

\begin{figure}[ht]
\centering
\begin{tikzpicture}[scale=0.95]
  % P2
  \begin{scope}[xshift=-4.3cm]
    \draw[->,black!55] (-1.55,0)--(2.55,0) node[below] {\(a\)};
    \draw[->,black!55] (0,-1.55)--(0,2.55) node[left] {\(b\)};
    \filldraw[fill=DictionaryBlue!13,draw=DictionaryBlue,very thick]
      (-1,-1)--(2,-1)--(-1,2)--cycle;
    \foreach \p in {(-1,-1),(2,-1),(-1,2)} {
      \fill \p circle (2.2pt);
    }
    \draw[fill=white,draw=black] (0,0) circle (2.3pt);
    \node[font=\bfseries] at (0.45,2.95) {\(\PP^2\)};
    \node[font=\small,align=center] at (0.45,-2.0)
      {three active facets\\normal fan unchanged};
  \end{scope}

  % P1xP1
  \begin{scope}[xshift=4.3cm]
    \draw[->,black!55] (-1.55,0)--(1.75,0) node[below] {\(a\)};
    \draw[->,black!55] (0,-1.55)--(0,1.75) node[left] {\(b\)};
    \filldraw[fill=PredictionTeal!13,draw=PredictionTeal,very thick]
      (-1,-1)--(1,-1)--(1,1)--(-1,1)--cycle;
    \foreach \p in {(-1,-1),(1,-1),(1,1),(-1,1)} {
      \fill \p circle (2.2pt);
    }
    \draw[fill=white,draw=black] (0,0) circle (2.3pt);
    \node[font=\bfseries] at (0.1,2.15)
      {\(\PP^1\times\PP^1\)};
    \node[font=\small,align=center] at (0.1,-2.0)
      {four active facets\\normal fan unchanged};
  \end{scope}
\end{tikzpicture}
\caption{The anticanonical polygons of \(\PP^2\) and
\(\PP^1\times\PP^1\). The open circle is the unique interior lattice
point. Every facet is at lattice distance one from the origin.}
\label{fig:first-anticanonical-polygons}
\end{figure}

\subsection{The Hirzebruch surfaces: ample, borderline, and negative}

Recall the fan of the Hirzebruch surface \(\mathbb F_r\), with rays in
cyclic order
\[
w_1=(1,0),
\qquad
w_2=(0,1),
\qquad
w_3=(-1,r),
\qquad
w_4=(0,-1).
\]
Its anticanonical polygon is
\[
\boxed{
P_r
=
\left\{
(x,y):
 x\geq-1,
\quad
-1\leq y\leq1,
\quad
x\leq1+ry
\right\}.
}
\]

The Cartier data on the four maximal cones are
\[
\begin{aligned}
m_{12}&=(-1,-1),
&
 m_{23}&=(1-r,-1),\\
m_{34}&=(r+1,1),
&
 m_{41}&=(-1,1).
\end{aligned}
\]
These four points tell the whole positivity story.

\begin{proposition}[Anticanonical positivity on \(\mathbb F_r\)]
\label{prop:hirzebruch-anticanonical-positivity}
For \(r\geq0\):
\[
\begin{array}{rcl}
\mathbb F_r\text{ is Fano}
&\Longleftrightarrow&
r=0\text{ or }1,\\[1mm]
-K_{\mathbb F_2}\text{ is basepoint free but not ample},&&\\[-1mm]
-K_{\mathbb F_r}\text{ is not basepoint free}
&\Longleftrightarrow&
r\geq3.
\end{array}
\]
\end{proposition}

\begin{proof}
For \(r=0\) and \(r=1\), all four Cartier data are distinct vertices
of \(P_r\), so the normal fan of the polygon is the original fan.
Thus \(-K\) is ample.

For \(r=2\),
\[
m_{12}=m_{23}=(-1,-1).
\]
All Cartier data still lie in \(P_2\), so \(-K\) is basepoint free,
but one wall is invisible and the support function is not strictly
convex. Hence \(-K\) is not ample.

For \(r\geq3\),
\[
m_{23}=(1-r,-1)
\]
fails the inequality \(x\geq-1\). Thus one local anticanonical frame
does not extend to a global section, and \(-K\) is not basepoint
free.
\end{proof}

\begin{figure}[ht]
\centering
\begin{tikzpicture}[scale=0.84]
  % F1
  \begin{scope}[xshift=-5.3cm,xscale=0.78]
    \draw[->,black!50] (-1.55,0)--(2.55,0) node[below] {\(x\)};
    \draw[->,black!50] (0,-1.45)--(0,1.65) node[left] {\(y\)};
    \filldraw[fill=DictionaryBlue!12,draw=DictionaryBlue,very thick]
      (-1,-1)--(0,-1)--(2,1)--(-1,1)--cycle;
    \foreach \p in {(-1,-1),(0,-1),(2,1),(-1,1)} {
      \fill \p circle (2.2pt);
    }
    \draw[fill=white,draw=black] (0,0) circle (2.2pt);
    \node[font=\bfseries] at (0.35,2.15) {\(\mathbb F_1\)};
    \node[font=\small,align=center] at (0.35,-1.95)
      {four vertices\\ample};
  \end{scope}

  % F2
  \begin{scope}[xscale=0.7]
    \draw[->,black!50] (-1.55,0)--(3.55,0) node[below] {\(x\)};
    \draw[->,black!50] (0,-1.45)--(0,1.65) node[left] {\(y\)};
    \filldraw[fill=PredictionTeal!12,draw=PredictionTeal,very thick]
      (-1,-1)--(3,1)--(-1,1)--cycle;
    \foreach \p in {(-1,-1),(3,1),(-1,1)} {
      \fill \p circle (2.2pt);
    }
    \draw[fill=white,draw=black] (0,0) circle (2.2pt);
    \draw[WarningRed,thick] (-1,-1) circle (5pt);
    \node[WarningRed,fill=white,inner sep=1pt,font=\scriptsize,anchor=south west]
      at (-0.82,-0.72) {two data};
    \node[font=\bfseries] at (0.65,2.15) {\(\mathbb F_2\)};
    \node[font=\small,align=center] at (0.65,-1.95)
      {one wall forgotten\\basepoint free, not ample};
  \end{scope}

  % F3
  \begin{scope}[xshift=5.3cm,xscale=0.65]
    \draw[->,black!50] (-2.45,0)--(4.55,0) node[below] {\(x\)};
    \draw[->,black!50] (0,-1.45)--(0,1.65) node[left] {\(y\)};
    \filldraw[fill=SurfaceGreen!11,draw=SurfaceGreen,very thick]
      (-1,-0.6667)--(4,1)--(-1,1)--cycle;
    \fill (-1,-0.6667) circle (2.2pt);
    \fill (4,1) circle (2.2pt);
    \fill (-1,1) circle (2.2pt);
    \draw[fill=white,draw=black] (0,0) circle (2.2pt);
    \fill[WarningRed] (-2,-1) circle (2.5pt);
    \draw[WarningRed,->,thick] (-1.45,-1.12)--(-1.88,-1.02);
    \node[WarningRed,fill=white,inner sep=1pt,font=\scriptsize,anchor=west]
      at (-1.42,-1.16) {outside datum};
    \node[font=\bfseries] at (0.8,2.15) {\(\mathbb F_3\)};
    \node[font=\small,align=center] at (0.8,-1.95)
      {a local frame does not globalize\\not basepoint free};
  \end{scope}
\end{tikzpicture}
\caption{The anticanonical polygons of \(\mathbb F_1\),
\(\mathbb F_2\), and \(\mathbb F_3\). At \(r=2\), two neighboring
Cartier data coincide. At \(r=3\), the red Cartier datum
\((-2,-1)\) lies outside the polytope.}
\label{fig:hirzebruch-anticanonical-trichotomy}
\end{figure}

There is also a curve-theoretic way to read the same transition. Let
\(S\) be the torus-invariant section with
\[
S^2=-r
\]
and let \(F\) be a fiber. Then
\[
-K_{\mathbb F_r}=2S+(r+2)F,
\]
so
\[
(-K_{\mathbb F_r})\cdot F=2,
\qquad
(-K_{\mathbb F_r})\cdot S=2-r.
\]
The anticanonical divisor is positive on \(S\) for \(r=0,1\), zero
on \(S\) for \(r=2\), and negative on \(S\) for \(r\geq3\). The
polygon and the invariant curve tell the same story.

\begin{surfacelaboratory}
The identification
\[
\mathbb F_1\simeq\Blp{\PP^2}{p}
\]
turns the first panel of \cref{fig:hirzebruch-anticanonical-trichotomy}
into the familiar corner-cut picture. On the blowup,
\[
-K_X=3H-E,
\]
and the representative used in Chapter~5 has polytope
\[
\left\{
(a,b):a,b\geq0,
\quad
1\leq a+b\leq3
\right\}.
\]
Translating by \((-1,-1)\) gives the centered anticanonical polygon.
Translation changes neither the normal fan nor the associated
projective map.
\end{surfacelaboratory}

\begin{pauseandthink}
The fan of \(\PP^2\) has three rays. Insert the ray \(e_1+e_2\) to
blow up one fixed point, and then insert a ray in a second maximal cone
to blow up another fixed point.

Write the new anticanonical inequalities. Does the second corner cut
still leave every inequality active? What does your answer predict
about the Fano property?
\end{pauseandthink}

The Fano condition has produced a very rigid kind of lattice polytope:
\[
\boxed{
0\text{ is the unique interior lattice point, and every facet is at
lattice distance }1\text{ from }0.
}
\]
The next section studies this condition intrinsically, without first
mentioning a toric variety.

\begin{prediction}
For each facet \(F\) of a Fano polytope \(P\subseteq M_\RR\), record
its primitive inward normal \(u_F\in N\), and form
\[
\Conv(u_F:F\text{ a facet of }P)\subseteq N_\RR.
\]

Compute this new polygon for the anticanonical triangle of \(\PP^2\)
and the square of \(\PP^1\times\PP^1\). What relation do you expect
between faces of the original polygon and faces of the new one? What
might happen if the construction is applied twice?
\end{prediction}
\section{Reflexive polytopes and polar duality}
\label{sec:reflexive-polytopes-polar-duality}
\index{polar duality}
\index{reflexive polytope}
\index{distance-one criterion}
\index{face duality}

The anticanonical polygons in the previous section all had the form
\[
P=
\left\{
 m\in M_\RR:
 \langle m,u_F\rangle\geq-1
 \text{ for every facet }F
\right\}.
\]
The primitive inward normals \(u_F\) were not merely auxiliary data:
they assembled into a second polytope in the dual vector space
\(N_\RR\). We now study that construction intrinsically.

\subsection{The polar remembers supporting hyperplanes}

Let \(P\subseteq M_\RR\) be a full-dimensional polytope with the
origin in its interior. Its \emph{polar polytope}, in the sign
convention used for toric divisors, is
\[
\boxed{
P^\circ
=
\left\{
 u\in N_\RR:
 \langle m,u\rangle\geq-1
 \text{ for every }m\in P
\right\}.
}
\]

The condition \(0\in\operatorname{int}(P)\) matters. It ensures that
\(P^\circ\) is bounded and full dimensional. It also puts the origin
in the interior of \(P^\circ\).

\begin{observation}
Some texts define the polar using
\[
\langle m,u\rangle\leq1.
\]
Replacing \(u\) by \(-u\) converts that convention into ours. The
present sign is convenient because toric divisor polytopes are already
written using inequalities
\[
\langle m,u_\rho\rangle\geq-a_\rho.
\]
\end{observation}

Suppose the facets of \(P\) have primitive inward normals \(u_F\in N\)
and facet presentation
\[
P=
\left\{
 m\in M_\RR:
 \langle m,u_F\rangle\geq-a_F
 \text{ for every facet }F
\right\},
\qquad a_F>0.
\]
Because \(u_F\) is primitive, the positive integer \(a_F\) is the
\emph{lattice distance} from the origin to the supporting hyperplane
of \(F\).

\begin{proposition}[Facets become vertices of the polar]
\label{prop:polar-from-facet-presentation}
With the notation above,
\[
\boxed{
P^\circ
=
\Conv
\left(
\frac{u_F}{a_F}:F\text{ a facet of }P
\right).
}
\]
In particular, the exact position of each facet determines a vertex of
the polar.
\end{proposition}

\noindent
\textit{Reference.}
This is the calculation in Cox--Little--Schenck, Exercise~2.2.1.

\begin{proof}
Set
\[
Q=
\Conv
\left(
\frac{u_F}{a_F}:F\text{ a facet of }P
\right).
\]
Then
\[
\begin{aligned}
Q^\circ
&=
\left\{
m\in M_\RR:
\left\langle m,\frac{u_F}{a_F}\right\rangle\geq-1
\text{ for every }F
\right\}\\
&=
\left\{
m\in M_\RR:
\langle m,u_F\rangle\geq-a_F
\text{ for every }F
\right\}\\
&=P.
\end{aligned}
\]
Polarity is involutive for full-dimensional polytopes containing the
origin in their interiors, so
\[
P^\circ=(Q^\circ)^\circ=Q.
\]
\end{proof}

The formula makes a crucial contrast visible:
\[
\boxed{
\begin{array}{c}
\text{The normal fan remembers the directions }u_F,\\[1mm]
\text{whereas the polar remembers the scaled vectors }u_F/a_F.
\end{array}
}
\]
Translations and positive dilations do not change a normal fan, but
they do change the polar.

\begin{pauseandthink}
Let \(P=[-1,2]\times[-1,1]\). Compute its polar from the four facet
inequalities. Which vertex of \(P^\circ\) fails to be a lattice point?

Now compare \(P\) with \([-1,1]^2\). The two rectangles have the same
normal fan. Why does polarity distinguish them?
\end{pauseandthink}

\subsection{The anticanonical triangle and its polar}

For \(X=\PP^2\), the anticanonical polygon is
\[
P=
\left\{
(x,y):x\geq-1,\quad y\geq-1,\quad x+y\leq1
\right\}.
\]
Its three primitive inward facet normals are
\[
u_1=(1,0),
\qquad
u_2=(0,1),
\qquad
u_0=(-1,-1).
\]
Every facet has lattice distance one from the origin, so
\[
P^\circ
=
\Conv(u_0,u_1,u_2).
\]
Thus the polar is exactly the convex hull of the ray generators of the
fan of \(\PP^2\).

\begin{figure}[ht]
\centering
\begin{tikzpicture}[scale=0.92]
  % P in M
  \begin{scope}[xshift=-4.5cm,scale=0.82]
    \draw[->,black!50] (-1.65,0)--(2.65,0) node[below] {\(x\)};
    \draw[->,black!50] (0,-1.65)--(0,2.65) node[left] {\(y\)};
    \fill[DictionaryBlue!8] (-1,-1)--(2,-1)--(-1,2)--cycle;

    \draw[DictionaryBlue,very thick] (-1,-1)--(-1,2);
    \draw[PredictionTeal,very thick] (-1,-1)--(2,-1);
    \draw[SurfaceGreen,very thick] (2,-1)--(-1,2);

    \foreach \p in {(-1,-1),(2,-1),(-1,2)} {
      \fill \p circle (2.3pt);
    }
    \draw[fill=white,draw=black] (0,0) circle (2.2pt);

    \node[DictionaryBlue,fill=white,inner sep=1pt] at (-1.36,0.55)
      {\(F_1\)};
    \node[PredictionTeal,fill=white,inner sep=1pt] at (0.55,-1.34)
      {\(F_2\)};
    \node[SurfaceGreen,fill=white,inner sep=1pt,rotate=-45] at (0.78,0.72)
      {\(F_0\)};
    \node[font=\bfseries] at (0.35,3.1) {\(P\subseteq M_\RR\)};
  \end{scope}

  \draw[->,very thick,black!55] (-0.9,0.15)--(0.9,0.15)
    node[midway,above,align=center,font=\small]
    {primitive\\facet normals};

  % polar in N
  \begin{scope}[xshift=4.5cm,scale=1.22]
    \draw[->,black!50] (-1.5,0)--(1.55,0) node[below] {\(e_1\)};
    \draw[->,black!50] (0,-1.5)--(0,1.55) node[left] {\(e_2\)};
    \filldraw[fill=ExampleGold!10,draw=black!55,thick]
      (1,0)--(0,1)--(-1,-1)--cycle;

    \fill[DictionaryBlue] (1,0) circle (2.5pt);
    \fill[PredictionTeal] (0,1) circle (2.5pt);
    \fill[SurfaceGreen] (-1,-1) circle (2.5pt);

    \node[DictionaryBlue,anchor=west,fill=white,inner sep=1pt]
      at (1.04,0.12) {\(u_1\)};
    \node[PredictionTeal,anchor=south west,fill=white,inner sep=1pt]
      at (0.06,1.02) {\(u_2\)};
    \node[SurfaceGreen,anchor=north east,fill=white,inner sep=1pt]
      at (-1.02,-1.02) {\(u_0\)};
    \draw[fill=white,draw=black] (0,0) circle (2.2pt);
    \node[font=\bfseries] at (0,2.08) {\(P^\circ\subseteq N_\RR\)};
  \end{scope}
\end{tikzpicture}
\caption{For the anticanonical triangle of \(\PP^2\), each colored
facet of \(P\) becomes the matching vertex of \(P^\circ\). The polar
is the convex hull of the primitive ray generators of the fan.}
\label{fig:p2-polar-facet-normal-duality}
\end{figure}

\subsection{Reflexive polytopes}

The previous example has a special integrality property: both \(P\)
and \(P^\circ\) are lattice polytopes.

\begin{definition}[Reflexive polytope]
\label{def:reflexive-polytope}
A full-dimensional lattice polytope \(P\subseteq M_\RR\) is
\emph{reflexive} if
\[
0\in\operatorname{int}(P)
\]
and its polar \(P^\circ\subseteq N_\RR\) is also a lattice polytope.
\end{definition}

\begin{theorem}[The distance-one criterion]
\label{thm:reflexive-distance-one-criterion}
Let \(P\subseteq M_\RR\) be a full-dimensional lattice polytope with
\(0\in\operatorname{int}(P)\). The following are equivalent:

\begin{enumerate}[label=(\alph*),wide=0pt]
\item \(P\) is reflexive.

\item Every facet of \(P\) is at lattice distance one from the origin.

\item The facet presentation of \(P\) is
\[
P=
\left\{
 m\in M_\RR:
 \langle m,u_F\rangle\geq-1
 \text{ for every facet }F
\right\},
\]
where every \(u_F\) is the primitive inward normal.
\end{enumerate}

When these conditions hold,
\[
\boxed{
P^\circ
=
\Conv(u_F:F\text{ a facet of }P).
}
\]
\end{theorem}

\noindent
\textit{Reference.}
This is Cox--Little--Schenck, Definition~2.3.12 together with
Exercise~2.2.1.

\begin{proof}
Write the facet presentation as
\[
P=
\{m:\langle m,u_F\rangle\geq-a_F\text{ for every }F\},
\]
with \(u_F\) primitive and \(a_F\in\ZZ_{>0}\). By
\cref{prop:polar-from-facet-presentation}, the vertices of the polar
are the vectors
\[
\frac{u_F}{a_F}.
\]
Since \(u_F\) is primitive, this vector lies in \(N\) if and only if
\(a_F=1\). Thus \(P^\circ\) is a lattice polytope exactly when every
facet has lattice distance one.
\end{proof}

\begin{corollary}[Reflexivity is symmetric]
\label{cor:polar-reflexive-involution}
If \(P\) is reflexive, then \(P^\circ\) is reflexive and
\[
(P^\circ)^\circ=P.
\]
Moreover, the origin is the unique interior lattice point of both
\(P\) and \(P^\circ\).
\end{corollary}

\begin{proof}
Both \(P\) and \(P^\circ\) are lattice polytopes containing the
origin in their interiors, and polarity is involutive. Hence the
definition is symmetric.

The unique-interior-point assertion follows from the distance-one
facet presentation exactly as in
\cref{prop:anticanonical-unique-interior-point}.
\end{proof}

\begin{warningbox}
Having the origin in the interior is not enough. For example,
\[
[-1,2]\times[-1,1]
\]
is a lattice polytope containing the origin, but the facet \(x=2\)
is at lattice distance two. Its polar has the nonintegral vertex
\[
\left(-\frac12,0\right),
\]
so the rectangle is not reflexive.
\end{warningbox}

\subsection{Faces reverse dimension}

Polarity exchanges faces of complementary dimensions. If
\(Q\subsetneq P\) is a nonempty face, define its \emph{dual face} by
\[
Q^*
=
\left\{
 u\in P^\circ:
 \langle m,u\rangle=-1
 \text{ for every }m\in Q
\right\}.
\]

\begin{proposition}[The polar face dictionary]
\label{prop:polar-face-dictionary}
Let \(P\) be reflexive of dimension \(n\). Then:

\begin{enumerate}[label=(\alph*),wide=0pt]
\item \(Q^*\) is a face of \(P^\circ\).

\item The correspondence \(Q\mapsto Q^*\) reverses inclusions.

\item
\[
\dim Q+\dim Q^*=n-1.
\]

\item Applying the construction twice returns the original face:
\[
(Q^*)^*=Q.
\]
\end{enumerate}
\end{proposition}

\begin{proof}[Conceptual proof sketch]
A face \(Q\) is cut out by the facets of \(P\) that contain it. The
primitive normals of precisely those facets are vertices of
\(P^\circ\), and their convex hull is \(Q^*\). Adding equations on
one side enlarges the dual face on the other, which reverses
inclusions. The dimension formula and involutivity are the usual dual
face statements for a polytope and its polar.
\end{proof}

For a polygon, the dictionary is especially visible:
\[
\boxed{
\text{edges of }P
\longleftrightarrow
\text{vertices of }P^\circ,
\qquad
\text{vertices of }P
\longleftrightarrow
\text{edges of }P^\circ.
}
\]

\begin{figure}[ht]
\centering
\begin{tikzpicture}[scale=1.0]
  % square
  \begin{scope}[xshift=-3.8cm,scale=1.15]
    \draw[->,black!45] (-1.5,0)--(1.55,0);
    \draw[->,black!45] (0,-1.5)--(0,1.55);
    \filldraw[fill=DictionaryBlue!8,draw=black!55,thick]
      (-1,-1)--(1,-1)--(1,1)--(-1,1)--cycle;
    \draw[DictionaryBlue,very thick] (1,-1)--(1,1);
    \fill[ExampleGold] (1,1) circle (3pt);
    \draw[fill=white,draw=black] (0,0) circle (2pt);
    \node[DictionaryBlue,anchor=west] at (1.08,0) {edge};
    \node[ExampleGold,anchor=south west] at (1.03,1.03) {vertex};
    \node[font=\bfseries] at (0,1.92) {\(P=[-1,1]^2\)};
  \end{scope}

  \draw[<->,very thick,black!55] (-0.85,0)--(0.85,0)
    node[midway,above,font=\small] {polar};

  % diamond
  \begin{scope}[xshift=3.8cm,scale=1.15]
    \draw[->,black!45] (-1.5,0)--(1.55,0);
    \draw[->,black!45] (0,-1.5)--(0,1.55);
    \filldraw[fill=PredictionTeal!8,draw=black!55,thick]
      (1,0)--(0,1)--(-1,0)--(0,-1)--cycle;
    \fill[DictionaryBlue] (-1,0) circle (3pt);
    \draw[ExampleGold,very thick] (-1,0)--(0,-1);
    \draw[fill=white,draw=black] (0,0) circle (2pt);
    \node[DictionaryBlue,anchor=east] at (-1.08,0) {vertex};
    \node[ExampleGold,anchor=north east,rotate=-45] at (-0.45,-0.56)
      {edge};
    \node[font=\bfseries] at (0,1.92) {\(P^\circ\)};
  \end{scope}
\end{tikzpicture}
\caption{For the square and diamond, the blue edge \(x=1\) becomes
the blue polar vertex \((-1,0)\). The gold vertex \((1,1)\) becomes
the gold edge joining \((-1,0)\) to \((0,-1)\).}
\label{fig:square-diamond-polar-face-duality}
\end{figure}

\subsection{The normal fan is already inside the polar}

The polar does more than record the normals one at a time. Its whole
face structure reconstructs the normal fan.

\begin{theorem}[Normal cones from polar faces]
\label{thm:normal-fan-cones-over-polar}
Let \(P\subseteq M_\RR\) be reflexive. For every face \(Q\subseteq P\),
the normal cone of \(Q\) is
\[
\boxed{
\sigma_Q=\Cone{Q^*}\subseteq N_\RR.
}
\]
Consequently, the normal fan \(\Sigma_P\) consists of the cones over
the proper faces of \(P^\circ\).
\end{theorem}

\noindent
\textit{Reference.}
This is Cox--Little--Schenck, Exercise~2.3.4.

\begin{proof}[Conceptual proof sketch]
A vector \(u\in N_\RR\) lies in \(\sigma_Q\) precisely when the
linear function \(\langle -,u\rangle\) reaches its minimum on \(P\)
along a face containing \(Q\). After rescaling a nonzero \(u\) until
that minimum is \(-1\), the resulting point lies in the dual face
\(Q^*\). Undoing the rescaling shows that the whole normal cone is the
cone over \(Q^*\).
\end{proof}

For the anticanonical polygon of a smooth toric Fano variety, this says
\[
\Sigma
=
\Sigma_{P_{-K_X}}
=
\{\text{cones over proper faces of }P_{-K_X}^\circ\}.
\]
In particular,
\[
P_{-K_X}^\circ
=
\Conv(u_\rho:\rho\in\Sigma(1)).
\]
The fan and the anticanonical polytope are therefore two sides of one
polar pair.

\begin{threedictionaries}
For a reflexive polytope, the three dictionaries become symmetric.

\medskip

\begin{center}
\small
\begin{tabularx}{\linewidth}{@{}
  >{\raggedright\arraybackslash}X
  >{\raggedright\arraybackslash}X
  >{\raggedright\arraybackslash}X@{}}
\toprule
\textbf{Polytope \(P\)}
&
\textbf{Polar \(P^\circ\)}
&
\textbf{Toric variety \(X_{\Sigma_P}\)}
\\
\midrule
facet \(F\)
&
vertex \(u_F\)
&
ray \(\rho_F\) and divisor \(D_{\rho_F}\)
\\
vertex \(v\)
&
facet \(v^*\)
&
maximal cone and torus fixed point
\\
face inclusion
&
reverse face inclusion
&
reverse orbit--cone correspondence
\\
lattice distance one
&
lattice vertex
&
integral anticanonical local data
\\
\bottomrule
\end{tabularx}
\end{center}
\end{threedictionaries}

\begin{prediction}
Suppose \(X_\Sigma\) is singular and its anticanonical inequalities
cut out a bounded polytope with normal fan \(\Sigma\). Its vertices
need not lie in \(M\).

On each maximal cone \(\sigma\), what condition would produce a lattice
vector \(m_\sigma\in M\) satisfying
\[
\langle m_\sigma,u_\rho\rangle=-1
\qquad
(\rho\in\sigma(1))?
\]
How is this the Cartier condition for \(-K_{X_\Sigma}\)?
\end{prediction}

\section{Gorenstein Fano toric varieties}
\label{sec:gorenstein-fano-toric-varieties}
\index{Gorenstein Fano toric variety}
\index{reflexive polytope!Gorenstein Fano correspondence}
\index{weighted projective space!Gorenstein criterion}

The previous two sections isolated two different features of the
anticanonical divisor:

\begin{itemize}
\item ampleness is detected by the strict convexity of the
anticanonical support function;

\item Cartierness is detected by integral height-one data on the
maximal cones.
\end{itemize}

A Gorenstein Fano toric variety is what we obtain when both conditions
hold at once. On the polytope side, the same combination becomes
reflexivity.

\subsection{Three levels of anticanonical positivity}

Let \(X\) be a complete normal variety.

\begin{definition}
The variety \(X\) is

\begin{enumerate}[label=(\alph*),wide=0pt]
\item \emph{\(\QGorenstein\)} if \(K_X\) is \(\QCartier\);

\item \emph{\(\QQ\)-Fano} if \(K_X\) is \(\QCartier\) and \(-K_X\)
is ample;

\item \emph{Gorenstein Fano} if \(-K_X\) is Cartier and ample.
\end{enumerate}

When \(X\) is smooth and \(-K_X\) is ample, we simply say that \(X\)
is Fano.
\end{definition}

\noindent
\textit{Reference.}
The Gorenstein Fano terminology is Cox--Little--Schenck,
Definition~8.3.1.

Thus
\[
\text{smooth Fano}
\Longrightarrow
\text{Gorenstein Fano}
\Longrightarrow
\QQ\text{-Fano}.
\]
Neither converse holds. The point of this section is to see the two
failures directly in a fan and its anticanonical polytope.

\subsection{Local height-one data become polytope vertices}

Let \(X_\Sigma\) be complete and \(\QGorenstein\). For each maximal
cone \(\sigma\), the height-one criterion from
\cref{prop:gorenstein-height-one} gives a rational vector
\[
h_\sigma\in M_\QQ
\]
such that
\[
\langle h_\sigma,u_\rho\rangle=1
\qquad
(\rho\in\sigma(1)).
\]
Equivalently, the local Cartier datum for \(-K_X\) is
\[
m_\sigma=-h_\sigma,
\qquad
\langle m_\sigma,u_\rho\rangle=-1
\qquad
(\rho\in\sigma(1)).
\]

If \(-K_X\) is ample, then \(m_\sigma\) is not merely a local linear
function: it is the vertex of \(P_{-K_X}\) whose normal cone is
\(\sigma\). Cartierness asks whether these vertices lie in the lattice
\(M\).

\begin{proposition}[Integrality of anticanonical vertices]
\label{prop:gorenstein-anticanonical-vertices}
Let \(X_\Sigma\) be a \(\QQ\)-Fano toric variety. The following are
equivalent:

\begin{enumerate}[label=(\alph*),wide=0pt]
\item \(X_\Sigma\) is Gorenstein Fano.

\item Every anticanonical Cartier datum \(m_\sigma\) belongs to \(M\).

\item Every vertex of \(P_{-K_X}\) is a lattice point.

\item For every maximal cone \(\sigma\), the primitive ray generators
of \(\sigma\) lie on an integral height-one hyperplane.
\end{enumerate}
\end{proposition}

\begin{proof}
The equivalence of \textup{(a)}, \textup{(b)}, and \textup{(d)} is
exactly the height-one criterion for the Cartier divisor \(-K_X\).
Since \(-K_X\) is ample, strict convexity identifies the local data
\(m_\sigma\) with the vertices of \(P_{-K_X}\). This gives the
equivalence with \textup{(c)}.
\end{proof}

This is the local-to-global mechanism behind reflexivity:
\[
\boxed{
\begin{array}{c}
\text{the equations }\langle m,u_\rho\rangle=-1
\text{ define rational vertices},\\[1mm]
\text{and Gorensteinness says that all of those vertices are integral.}
\end{array}
}
\]

\subsection{Gorenstein Fano varieties and reflexive polytopes}

We can now state the central theorem of the chapter.

\begin{theorem}[Gorenstein Fano--reflexive correspondence]
\label{thm:gorenstein-fano-reflexive-correspondence}
The following constructions are inverse to one another, up to lattice
equivalence.

\begin{enumerate}[label=(\alph*),wide=0pt]
\item If \(X_\Sigma\) is a projective Gorenstein Fano toric variety,
then its anticanonical polytope
\[
P_{-K_X}
=
\left\{
 m\in M_\RR:
 \langle m,u_\rho\rangle\geq-1
 \text{ for every }\rho
\right\}
\]
is reflexive and has normal fan \(\Sigma\).

\item If \(P\subseteq M_\RR\) is reflexive, then the projective toric
variety \(X_{\Sigma_P}\) is Gorenstein Fano, and
\[
P=P_{-K_{X_{\Sigma_P}}}.
\]
\end{enumerate}
\end{theorem}

\noindent
\textit{Reference.}
This is Cox--Little--Schenck, Theorem~8.3.4.

\begin{proof}
Suppose first that \(X_\Sigma\) is Gorenstein Fano. Since \(-K_X\) is
Cartier, its polytope is a lattice polytope. Since it is ample, every
ray inequality is active and the normal fan of the polytope is exactly
\(\Sigma\). Its facet presentation is
\[
P_{-K_X}
=
\left\{
 m\in M_\RR:
 \langle m,u_\rho\rangle\geq-1
 \text{ for every }\rho
\right\},
\]
with each \(u_\rho\) primitive. Every facet is therefore at lattice
distance one from the origin. By
\cref{thm:reflexive-distance-one-criterion}, the polytope is reflexive.

Conversely, let \(P\) be reflexive. Its facet presentation is
\[
P=
\left\{
 m\in M_\RR:
 \langle m,u_F\rangle\geq-1
 \text{ for every facet }F
\right\},
\]
where \(u_F\in N\) is the primitive inward normal. The divisor attached
to this polytope on \(X_{\Sigma_P}\) is therefore
\[
D_P=\sum_F D_F=-K_{X_{\Sigma_P}}.
\]
Because \(P\) is a lattice polytope, this divisor is Cartier; because
\(\Sigma_P\) is its normal fan, it is ample. Hence
\(X_{\Sigma_P}\) is Gorenstein Fano.
\end{proof}

For a Gorenstein Fano toric variety, polarity takes the particularly
concrete form
\[
\boxed{
P_{-K_X}^\circ
=
\Conv\bigl(u_\rho:\rho\in\Sigma(1)\bigr).
}
\]
Moreover,
\[
\Sigma
=
\{\text{cones over the proper faces of }P_{-K_X}^\circ\}.
\]
Thus the polar simultaneously records the ray generators and how they
assemble into cones.

\begin{observation}
Since the polar of a reflexive polytope is again reflexive, a reflexive
pair
\[
P\subseteq M_\RR,
\qquad
P^\circ\subseteq N_\RR
\]
produces two Gorenstein Fano toric varieties. They need not be
isomorphic. This symmetry is one reason reflexive polytopes play a
central role in toric mirror symmetry; no familiarity with mirror
symmetry is needed here.
\end{observation}

\begin{threedictionaries}
The Gorenstein Fano condition is the same phenomenon in three
languages.

\medskip

\begin{center}
\small
\begin{tabularx}{\linewidth}{@{}
  >{\raggedright\arraybackslash}X
  >{\raggedright\arraybackslash}X
  >{\raggedright\arraybackslash}X@{}}
\toprule
\textbf{Algebra}
&
\textbf{Geometry}
&
\textbf{Combinatorics}
\\
\midrule
\(-K_X\) is Cartier
&
\(X\) is Gorenstein
&
local height-one data are integral
\\
\(-K_X\) is ample
&
anticanonical powers embed \(X\)
&
its support function is strictly convex
\\
Cartier and ample together
&
\(X\) is Gorenstein Fano
&
\(P_{-K_X}\) is reflexive
\\
local anticanonical frames
&
torus fixed points
&
vertices of \(P_{-K_X}\)
\\
\bottomrule
\end{tabularx}
\end{center}
\end{threedictionaries}

\subsection{Weighted projective spaces}

Weighted projective spaces show that ampleness and Cartierness are
independent questions. Let
\[
X=\PP(q_0,\dots,q_n),
\qquad
\gcd(q_0,\dots,q_n)=1.
\]
Its Cox ring is
\[
S=\Bbbk[x_0,\dots,x_n],
\qquad
\deg(x_i)=q_i\in\Cl(X)\simeq\ZZ.
\]
Since the coordinate divisor \(D_i\) has class \(q_i\),
\[
[-K_X]
=
\left[\sum_iD_i\right]
=
(q_0+\cdots+q_n)\in\Cl(X).
\]
This class is positive, so it is ample as a rational divisor. The
remaining question is whether it is Cartier.

\begin{proposition}[The weighted-projective Gorenstein criterion]
\label{prop:weighted-projective-gorenstein-fano}
The weighted projective space \(\PP(q_0,\dots,q_n)\) is Gorenstein
Fano if and only if
\[
q_i\mid q_0+\cdots+q_n
\qquad
\text{for every }i.
\]
Every weighted projective space is \(\QQ\)-Fano.
\end{proposition}

\noindent
\textit{Reference.}
This is Cox--Little--Schenck, Example~8.3.3 and Exercise~8.3.2.

\begin{proof}[Conceptual proof sketch]
The anticanonical class is the positive class
\[
q_0+\cdots+q_n,
\]
so it is ample. A class of degree \(d\) defines a line bundle on the
weighted projective space precisely when every weight divides \(d\),
or equivalently when
\[
\operatorname{lcm}(q_0,\dots,q_n)\mid d.
\]
Applying this to \(d=q_0+\cdots+q_n\) gives the criterion.
\end{proof}

\subsection{A singular comparison: \texorpdfstring{\(\PP(1,1,2)\)}{P(1,1,2)} and
\texorpdfstring{\(\PP(1,1,3)\)}{P(1,1,3)}}

For \(X_n=\PP(1,1,n)\), choose the ray generators
\[
u_0=(1,0),
\qquad
u_1=(-1,n),
\qquad
u_2=(0,-1).
\]
They satisfy
\[
u_0+u_1+n u_2=0.
\]
The anticanonical region is
\[
P_{-K_{X_n}}
=
\left\{
(x,y):
 x\geq-1,
\quad
-x+ny\geq-1,
\quad
y\leq1
\right\}.
\]
Its vertices are
\[
(-1,1),
\qquad
(n+1,1),
\qquad
\left(-1,-\frac{2}{n}\right).
\]
Thus
\[
P_{-K_{X_n}}
\text{ is a lattice polygon}
\quad\Longleftrightarrow\quad
n\mid2.
\]
This recovers the conclusion from Chapter~4:
\[
\PP(1,1,2)
\text{ is singular Gorenstein Fano},
\]
whereas
\[
\PP(1,1,3)
\text{ is \(\QQ\)-Fano but not Gorenstein}.
\]

\begin{figure}[ht]
\centering
\begin{tikzpicture}[scale=0.9]
  % P(1,1,2)
  \begin{scope}[xshift=-4.4cm,scale=0.72]
    \draw[step=1,black!10] (-1.6,-1.6) grid (3.8,1.7);
    \draw[->,black!50] (-1.7,0)--(4.0,0) node[below] {\(x\)};
    \draw[->,black!50] (0,-1.7)--(0,1.8) node[left] {\(y\)};

    \filldraw[fill=PredictionTeal!12,draw=PredictionTeal,very thick]
      (-1,-1)--(-1,1)--(3,1)--cycle;
    \foreach \p in {(-1,-1),(-1,1),(3,1)} {
      \fill \p circle (3pt);
    }
    \draw[fill=white,draw=black] (0,0) circle (2.5pt);

    \node[font=\bfseries] at (0.7,2.35) {\(\PP(1,1,2)\)};
    \node[PredictionTeal,align=center] at (1.0,-1.48)
      {all vertices are integral};
  \end{scope}

  % arrow
  \draw[->,very thick,black!55] (-0.8,0.1)--(0.8,0.1)
    node[midway,above,font=\small,align=center]
    {change the\\third weight};

  % P(1,1,3)
  \begin{scope}[xshift=4.4cm,scale=0.72]
    \draw[step=1,black!10] (-1.6,-1.6) grid (4.8,1.7);
    \draw[->,black!50] (-1.7,0)--(5.0,0) node[below] {\(x\)};
    \draw[->,black!50] (0,-1.7)--(0,1.8) node[left] {\(y\)};

    \filldraw[fill=ExampleGold!12,draw=ExampleGold!80!black,very thick]
      (-1,-0.6667)--(-1,1)--(4,1)--cycle;
    \fill (-1,1) circle (3pt);
    \fill (4,1) circle (3pt);
    \draw[fill=white,draw=WarningRed,very thick]
      (-1,-0.6667) circle (3.5pt);
    \draw[fill=white,draw=black] (0,0) circle (2.5pt);

    \node[WarningRed,anchor=east,fill=white,inner sep=1pt]
      at (-1.18,-0.58) {\(\left(-1,-\frac23\right)\)};
    \draw[WarningRed,->,thick] (-1.12,-0.60)--(-1.02,-0.66);
    \node[font=\bfseries] at (1.1,2.35) {\(\PP(1,1,3)\)};
    \node[ExampleGold!80!black,align=center] at (1.55,-1.48)
      {one rational vertex};
  \end{scope}
\end{tikzpicture}
\caption{The anticanonical polygons of two singular weighted
projective planes. The polygon for \(\PP(1,1,2)\) is reflexive. For
\(\PP(1,1,3)\), ampleness survives but one anticanonical vertex leaves
the lattice, recording the failure of Gorensteinness.}
\label{fig:weighted-projective-gorenstein-comparison}
\end{figure}

\begin{pauseandthink}
The singular cone of \(\PP(1,1,n)\) is generated by
\[
(1,0)
\qquad\text{and}\qquad
(-1,n).
\]
Solve the height-one equations on this cone and compare the resulting
functional with the nonintegral vertex
\[
\left(-1,-\frac{2}{n}\right)
\]
of the anticanonical polytope. Why do the two calculations contain the
same divisibility condition?
\end{pauseandthink}

\subsection{Reflexive does not mean smooth}

The polygon for \(\PP(1,1,2)\) is reflexive, but the corresponding
surface is singular. Reflexivity guarantees that \(-K_X\) is Cartier
and ample; it does not say that the cones of the fan are smooth.

Let \(P\subseteq M_\RR\) be reflexive. Since
\[
\Sigma_P
=
\{\text{cones over faces of }P^\circ\},
\]
the toric variety \(X_{\Sigma_P}\) is smooth precisely when every
maximal cone over a facet of \(P^\circ\) is smooth. Equivalently,

\[
\boxed{
X_{\Sigma_P}\text{ is smooth}
\quad\Longleftrightarrow\quad
\text{every facet of }P^\circ
\text{ is a unimodular simplex.}
}
\]

For a polygon, this becomes especially concrete. If
\[
u_1,\dots,u_s
\]
are the vertices of \(P^\circ\) in cyclic order, then
\[
X_{\Sigma_P}\text{ is smooth}
\quad\Longleftrightarrow\quad
|\det(u_i,u_{i+1})|=1
\qquad
\text{for every }i.
\]

Thus a smooth toric Fano surface corresponds not merely to a reflexive
polygon, but to a reflexive polygon whose polar has unimodular edges.
That is the starting point for the final surface laboratory.

\begin{prediction}
Begin with the anticanonical triangle of \(\PP^2\). A toric blowup cuts
off one corner and inserts the sum of the two neighboring ray
generators.

Which torus-fixed corner cuts preserve both of the following
properties?
\[
\begin{aligned}
&\text{every facet remains at lattice distance one from the origin},\\
&\text{every consecutive pair of polar vertices remains a lattice basis}.
\end{aligned}
\]
How many times can this be done before the anticanonical divisor stops
being ample?
\end{prediction}

\section{Toric del Pezzo surfaces through their polygons}
\label{sec:toric-del-pezzo-polygons}
\index{del Pezzo surface!toric}
\index{degree!del Pezzo surface}
\index{lattice polygon!area and degree}
\index{corner cut!anticanonical polygon}

A \emph{del Pezzo surface} is a smooth projective surface whose
anticanonical divisor is ample. For a smooth toric surface, the
previous sections turn this definition into a problem about a lattice
polygon:
\[
X_\Sigma\text{ is del Pezzo}
\quad\Longleftrightarrow\quad
P_{-K_X}\text{ is reflexive and }\Sigma_{P_{-K_X}}=\Sigma.
\]
Smoothness asks that consecutive primitive ray generators form lattice
bases. Ampleness asks that every wall of the fan remain visible to the
anticanonical polygon. In dimension two, these two restrictions leave
only five possibilities.

\subsection{The five smooth toric del Pezzo surfaces}

\begin{theorem}[Classification of smooth toric del Pezzo surfaces]
\label{thm:classification-smooth-toric-del-pezzo}
Up to isomorphism, the smooth toric del Pezzo surfaces are
\[
\boxed{
\PP^2,
\qquad
\PP^1\times\PP^1,
\qquad
\Blp{\PP^2}{p},
\qquad
\Blp{\PP^2}{p,q},
\qquad
\Blp{\PP^2}{p,q,r},
}
\]
where \(p,q,r\) are distinct torus-fixed points of \(\PP^2\), with
as many of them blown up as the notation indicates.
Their anticanonical degrees are respectively
\[
9,\qquad 8,\qquad 8,\qquad 7,\qquad 6.
\]
\end{theorem}

\noindent
\textit{Reference.}
The polytope classification is Cox--Little--Schenck,
Exercises~8.3.6--8.3.8. In particular, Exercise~8.3.7 isolates the five
smooth cases among the sixteen reflexive polygons, and
Exercise~8.3.8 identifies the corresponding toric surfaces.

\begin{proof}[Conceptual proof sketch]
A smooth toric Fano surface corresponds, by
\cref{thm:gorenstein-fano-reflexive-correspondence}, to a reflexive
polygon whose polar has unimodular edges. The classification of
reflexive polygons leaves exactly five such lattice-equivalence
classes.

There is also a geometric way to organize the same list. Every smooth
complete toric surface is obtained from \(\PP^2\), \(\PP^1\times\PP^1\),
or a Hirzebruch surface by toric blowups; this is
Cox--Little--Schenck, Theorem~10.4.3. By
\cref{prop:hirzebruch-anticanonical-positivity}, the only Fano
Hirzebruch surfaces are
\[
\mathbb F_0\simeq\PP^1\times\PP^1
\qquad\text{and}\qquad
\mathbb F_1\simeq\Blp{\PP^2}{p}.
\]
Starting from \(\PP^2\), at most three successive torus-fixed blowups
preserve ampleness of the anticanonical divisor. These produce the
remaining three surfaces in the list.
\end{proof}

\begin{figure}[ht]
\centering
\begin{tikzpicture}[scale=0.67]
  % P2
  \begin{scope}[xshift=-7.2cm,yshift=3.25cm]
    \filldraw[fill=DictionaryBlue!13,draw=DictionaryBlue,very thick]
      (-1,-1)--(2,-1)--(-1,2)--cycle;
    \foreach \p in {(-1,-1),(2,-1),(-1,2)} {\fill \p circle (2.2pt);}
    \draw[fill=white,draw=black] (0,0) circle (2.2pt);
    \node[font=\bfseries] at (0.35,2.65) {\(\PP^2\)};
    \node[font=\small] at (0.35,-1.55) {\((-K)^2=9\)};
  \end{scope}

  % P1xP1
  \begin{scope}[xshift=-1.8cm,yshift=3.25cm]
    \filldraw[fill=PredictionTeal!13,draw=PredictionTeal,very thick]
      (-1,-1)--(1,-1)--(1,1)--(-1,1)--cycle;
    \foreach \p in {(-1,-1),(1,-1),(1,1),(-1,1)} {\fill \p circle (2.2pt);}
    \draw[fill=white,draw=black] (0,0) circle (2.2pt);
    \node[font=\bfseries] at (0,1.65) {\(\PP^1\times\PP^1\)};
    \node[font=\small] at (0,-1.55) {\((-K)^2=8\)};
  \end{scope}

  % Bl1
  \begin{scope}[xshift=4.2cm,yshift=3.25cm]
    \filldraw[fill=SurfaceGreen!13,draw=SurfaceGreen,very thick]
      (-1,0)--(0,-1)--(2,-1)--(-1,2)--cycle;
    \foreach \p in {(-1,0),(0,-1),(2,-1),(-1,2)} {\fill \p circle (2.2pt);}
    \draw[fill=white,draw=black] (0,0) circle (2.2pt);
    \node[font=\bfseries] at (0.4,2.65) {\(\Blp{\PP^2}{p}\)};
    \node[font=\small] at (0.4,-1.55) {\((-K)^2=8\)};
  \end{scope}

  % Bl2
  \begin{scope}[xshift=-4.2cm,yshift=-2.6cm]
    \filldraw[fill=ExampleGold!18,draw=ExampleGold!80!black,very thick]
      (-1,0)--(0,-1)--(1,-1)--(1,0)--(-1,2)--cycle;
    \foreach \p in {(-1,0),(0,-1),(1,-1),(1,0),(-1,2)} {\fill \p circle (2.2pt);}
    \draw[fill=white,draw=black] (0,0) circle (2.2pt);
    \node[font=\bfseries] at (0,2.65) {\(\Blp{\PP^2}{p,q}\)};
    \node[font=\small] at (0,-1.55) {\((-K)^2=7\)};
  \end{scope}

  % Bl3
  \begin{scope}[xshift=3.3cm,yshift=-2.6cm]
    \filldraw[fill=PausePurple!12,draw=PausePurple,very thick]
      (-1,0)--(-1,1)--(0,1)--(1,0)--(1,-1)--(0,-1)--cycle;
    \foreach \p in {(-1,0),(-1,1),(0,1),(1,0),(1,-1),(0,-1)} {\fill \p circle (2.2pt);}
    \draw[fill=white,draw=black] (0,0) circle (2.2pt);
    \node[font=\bfseries] at (0,1.65) {\(\Blp{\PP^2}{p,q,r}\)};
    \node[font=\small] at (0,-1.55) {\((-K)^2=6\)};
  \end{scope}
\end{tikzpicture}
\caption{The five anticanonical polygons of smooth toric del Pezzo
surfaces, up to lattice equivalence. The open circle is the origin,
the unique interior lattice point.}
\label{fig:five-smooth-toric-del-pezzo-polygons}
\end{figure}

The two degree-eight surfaces are genuinely different. The polygon of
\(\PP^1\times\PP^1\) is centrally symmetric, while the polygon of
\(\Bl_p(\PP^2)\simeq\mathbb F_1\) is not. Equivalently, the first
surface has no curve of self-intersection \(-1\), while the second
contains the exceptional curve.

\subsection{Degree, area, and boundary lattice points}

For any ample divisor on a projective toric surface,
\cref{thm:polygon-area-self-intersection} gives
\[
D^2=2\operatorname{Area}(P_D).
\]
For an anticanonical reflexive polygon, the origin is the unique
interior lattice point. Pick's theorem therefore gives
\[
\operatorname{Area}(P_{-K_X})
=
1+\frac{1}{2}|\partial P_{-K_X}\cap M|-1
=
\frac{1}{2}|\partial P_{-K_X}\cap M|.
\]
Hence
\[
\boxed{
(-K_X)^2
=
|\partial P_{-K_X}\cap M|.
}
\]
The degree of a smooth toric del Pezzo surface is literally the number
of lattice points around the boundary of its anticanonical polygon.
Since the only interior lattice point is the origin,
\[
h^0(X,-K_X)
=
|P_{-K_X}\cap M|
=
(-K_X)^2+1.
\]

\begin{pauseandthink}
Check the degree formula directly on the five polygons in
\cref{fig:five-smooth-toric-del-pezzo-polygons}. Remember that a lattice
edge of length \(\ell\) contributes \(\ell\) to the number of boundary
lattice points when the edge lengths are summed cyclically.
\end{pauseandthink}

There is also a curve-by-curve version. By
\cref{thm:edge-length-intersection}, if the facet \(F_\rho\) of
\(P_{-K_X}\) corresponds to the invariant curve \(D_\rho\), then
\[
(-K_X)\cdot D_\rho
=
\length_M(F_\rho).
\]
Thus ampleness of \(-K_X\) is visible in the polygon as the positivity
of every edge length.

\subsection{Toric blowups are anticanonical corner cuts}

Suppose a fixed point corresponds to a smooth cone
\[
\sigma=\Cone{u,v},
\qquad
\det(u,v)=1.
\]
The toric blowup inserts the ray generated by \(u+v\). For the
anticanonical polygon, this adds the inequality
\[
\langle m,u+v\rangle\geq-1.
\]
At the old corner, both
\[
\langle m,u\rangle=-1
\qquad\text{and}\qquad
\langle m,v\rangle=-1,
\]
so the new inequality cuts that corner off. In the smooth case, the
removed triangle is unimodular and has area \(1/2\). Consequently,
\[
(-K_{\widetilde X})^2
=
(-K_X)^2-1,
\]
as expected from the blowup formula
\[
K_{\widetilde X}=\pi^*K_X+E.
\]

\begin{figure}[ht]
\centering
\begin{tikzpicture}[scale=0.72]
  \def\sep{4.25}
  % P2
  \begin{scope}[xshift=-6.35cm]
    \filldraw[fill=DictionaryBlue!12,draw=DictionaryBlue,very thick]
      (-1,-1)--(2,-1)--(-1,2)--cycle;
    \draw[WarningRed,dashed,very thick] (-1,0)--(0,-1);
    \node[font=\bfseries] at (0.35,2.55) {degree \(9\)};
    \node[font=\scriptsize,align=center] at (0.35,-1.62) {first\\corner cut};
  \end{scope}

  % Bl1
  \begin{scope}[xshift=-2.1cm]
    \filldraw[fill=SurfaceGreen!12,draw=SurfaceGreen,very thick]
      (-1,0)--(0,-1)--(2,-1)--(-1,2)--cycle;
    \draw[WarningRed,dashed,very thick] (1,-1)--(1,0);
    \node[font=\bfseries] at (0.35,2.55) {degree \(8\)};
    \node[font=\scriptsize,align=center] at (0.35,-1.62) {second\\corner cut};
  \end{scope}

  % Bl2
  \begin{scope}[xshift=2.15cm]
    \filldraw[fill=ExampleGold!16,draw=ExampleGold!80!black,very thick]
      (-1,0)--(0,-1)--(1,-1)--(1,0)--(-1,2)--cycle;
    \draw[WarningRed,dashed,very thick] (-1,1)--(0,1);
    \node[font=\bfseries] at (0,2.55) {degree \(7\)};
    \node[font=\scriptsize,align=center] at (0,-1.62) {third\\corner cut};
  \end{scope}

  % Bl3
  \begin{scope}[xshift=6.35cm]
    \filldraw[fill=PausePurple!11,draw=PausePurple,very thick]
      (-1,0)--(-1,1)--(0,1)--(1,0)--(1,-1)--(0,-1)--cycle;
    \draw[WarningRed,dashed,very thick] (-1,1)--(0,-1);
    \node[font=\bfseries] at (0,2.55) {degree \(6\)};
    \node[font=\scriptsize,align=center] at (0,-1.72)
      {a fourth cut collapses\\two adjacent facets};
  \end{scope}

  \foreach \x in {-4.25,0,4.25} {
    \draw[->,black!55,thick] (\x-0.45,0)--(\x+0.45,0);
  }
\end{tikzpicture}
\caption{Successive torus-fixed blowups of \(\PP^2\) cut unimodular
corners from the anticanonical polygon and lower the degree by one. At
degree six, every edge has lattice length one; a further blowup makes
two neighboring old facets collapse, so the new anticanonical divisor
is no longer ample.}
\label{fig:del-pezzo-anticanonical-corner-cuts}
\end{figure}

Why does the process stop after the third blowup? On
\(\Bl_{p,q,r}(\PP^2)\), every edge of the anticanonical hexagon has
lattice length one. Blow up the fixed point where two such invariant
curves meet. Their strict transforms satisfy
\[
(-K_{\widetilde X})\cdot D'_i
=
(-K_X)\cdot D_i-1
=
0.
\]
Therefore \(-K_{\widetilde X}\) cannot be ample. In the polygon, the
same failure appears when two old facets shrink to length zero. The
new support function remains convex, but it is no longer strictly
convex along those walls.

\begin{surfacelaboratory}
For a smooth toric del Pezzo surface, one polygon simultaneously shows:
\[
\begin{array}{c@{\qquad}c}
\text{facet normals} & \text{the fan},\\
\text{edge lengths} & (-K_X)\cdot D_\rho,\\
\text{twice the area} & (-K_X)^2,\\
\text{lattice points} & H^0(X,-K_X),\\
\text{corner cuts} & \text{toric blowups}.
\end{array}
\]
This is the surface form of all three dictionaries at once.
\end{surfacelaboratory}

\begin{prediction}
For each of the five polygons \(P=P_{-K_X}\), draw the polar polygon
\(P^\circ\) and count boundary lattice points on both sides. The five
examples suggest the identity
\[
|\partial P\cap M|
+
|\partial P^\circ\cap N|
=12.
\]
Why should the number \(12\) appear? We will return to this question
when Riemann--Roch turns lattice-point counts into intersection
numbers.
\end{prediction}

\begin{guidingquestion}
The anticanonical polygon controls global sections and positivity, but
what can it tell us about the topology of the underlying toric
variety---its fixed points, Betti numbers, and cohomology ring?
\end{guidingquestion}

\clearpage
\section*{What this chapter has accomplished}
\addcontentsline{toc}{section}{What this chapter has accomplished}
\markboth{What this chapter has accomplished}{What this chapter has accomplished}

The toric boundary is not merely a collection of invariant divisors.
It is the divisor along which invariant differential forms acquire
logarithmic poles. For every character \(\character{m}\),
\[
\frac{d\character{m}}{\character{m}}
\]
extends as a logarithmic one-form, and its residue along \(D_\rho\) is
\[
\operatorname{Res}_{D_\rho}
\left(\frac{d\character{m}}{\character{m}}\right)
=
\langle m,u_\rho\rangle.
\]
On a smooth toric variety, these invariant forms trivialize the
logarithmic cotangent bundle:
\[
\Omega^1_{X_\Sigma}(\log\boundaryX)
\simeq
M\otimes_{\ZZ}\mathcal O_{X_\Sigma}.
\]
Taking the top exterior power gives the conceptual form of the
canonical-divisor formula:
\[
\boxed{K_{X_\Sigma}+\boundaryX\sim0.}
\]

The anticanonical divisor therefore has the particularly uniform
polytope
\[
P_{-K_X}
=
\left\{
 m\in M_\RR:
 \langle m,u_\rho\rangle\geq-1
 \text{ for every }\rho
\right\}.
\]
Its positivity is read from its support function:
\[
X_\Sigma\text{ Fano}
\quad\Longleftrightarrow\quad
-K_{X_\Sigma}\text{ ample}
\quad\Longleftrightarrow\quad
\Sigma_{P_{-K_X}}=\Sigma.
\]
The origin is the unique interior lattice point of this polytope.

Polarity remembers more than the normal fan. If
\[
P=\{m:\langle m,u_F\rangle\geq-a_F\},
\]
then
\[
P^\circ
=
\Conv\left(\frac{u_F}{a_F}\right).
\]
Thus the normal fan remembers only the directions of the facet
normals, while the polar remembers their lattice distances from the
origin. A lattice polytope is reflexive precisely when every facet is
at lattice distance one from the origin.

This gives the central correspondence of the chapter:
\[
\boxed{
\begin{aligned}
P\text{ reflexive}
&\Longleftrightarrow
X_{\Sigma_P}\text{ Gorenstein Fano},\\
P^\circ\text{ lattice}
&\Longleftrightarrow
K_{X_{\Sigma_P}}\text{ Cartier}.
\end{aligned}}
\]
Smoothness is stronger: the relevant cones must be unimodular, not
merely Gorenstein.

In dimension two, the dictionary becomes completely visible. The five
smooth toric del Pezzo surfaces are
\[
\PP^2,
\qquad
\PP^1\times\PP^1,
\qquad
\Blp{\PP^2}{p},
\qquad
\Blp{\PP^2}{p,q},
\qquad
\Blp{\PP^2}{p,q,r},
\]
with anticanonical degrees
\[
9,\qquad8,\qquad8,\qquad7,\qquad6.
\]
Their anticanonical polygons record the fan, the invariant
intersection numbers, the global sections, and the toric blowups all
at once.

\begin{exercises}[Comprehension exercises]
\markboth{Comprehension exercises}{Comprehension exercises}

\item On the torus \(T=(\Gm)^2\), let
\[
\chi^{(a,b)}=x^ay^b.
\]
Compute
\[
\frac{d\chi^{(a,b)}}{\chi^{(a,b)}}.
\]
Compactify \(T\) to \(\PP^1\times\PP^1\), and compute the residue of
this logarithmic form along each of the four boundary divisors.
Explain how the four residues recover the pairings with the four ray
generators.

\item Let \(X_\Sigma\) be a smooth complete toric variety of dimension
\(n\). Choose a basis \(m_1,\dots,m_n\) of \(M\), and set
\[
\omega
=
\frac{d\character{m_1}}{\character{m_1}}
\wedge\cdots\wedge
\frac{d\character{m_n}}{\character{m_n}}.
\]
Show on a smooth affine chart that \(\omega\) has a simple pole along
each boundary component and no other zero or pole. Recover
\[
K_{X_\Sigma}=-\sum_\rho D_\rho.
\]

\item For \(X=\PP^2\), compute \(P_{-K_X}\) directly from the three
ray generators. Then:
\begin{enumerate}[label=(\alph*),wide=0pt]
\item list all lattice points of \(P_{-K_X}\);
\item verify that the origin is its unique interior lattice point;
\item recover \((-K_X)^2\) from its area;
\item recover \(h^0(X,-K_X)\) from its lattice points.
\end{enumerate}

\item For the Hirzebruch surface \(\mathbb F_r\), use the rays
\[
(1,0),\quad(0,1),\quad(-1,r),\quad(0,-1)
\]
to rederive
\[
P_{-K_{\mathbb F_r}}
=
\{(x,y):x\geq-1,\ -1\leq y\leq1,\ x\leq1+ry\}.
\]
Draw the cases \(r=0,1,2,3\). Determine directly from the polygons
when \(-K\) is ample, basepoint free but not ample, or not basepoint
free. Compare with
\[
(-K_{\mathbb F_r})\cdot S=2-r.
\]

\item Let
\[
P=\{m\in M_\RR:\langle m,u_i\rangle\geq-a_i,
\ i=1,\dots,s\}
\]
be full dimensional with \(0\in\operatorname{int}(P)\), where the
inequalities are irredundant and the \(u_i\) are primitive. Prove
\[
P^\circ
=
\Conv\left(\frac{u_1}{a_1},\dots,
\frac{u_s}{a_s}\right).
\]
Deduce the lattice-distance-one criterion for reflexivity.

\item Let \(P\subseteq M_\RR\) be reflexive. For a face \(Q\preceq P\),
define
\[
Q^*
=
\{u\in P^\circ:\langle m,u\rangle=-1
\text{ for every }m\in Q\}.
\]
Show in dimension two that vertices of \(P\) correspond to edges of
\(P^\circ\), and edges of \(P\) correspond to vertices of
\(P^\circ\). Verify
\[
\dim Q+\dim Q^*=1.
\]
Work this out explicitly for the square and diamond.

\item Consider the weighted projective plane \(\PP(1,1,n)\).
\begin{enumerate}[label=(\alph*),wide=0pt]
\item Show that it is \(\QQ\)-Fano for every \(n\geq1\).
\item Use the divisibility criterion to determine for which \(n\) it
is Gorenstein Fano.
\item Compute the anticanonical polygon for \(n=2\) and \(n=3\), and
locate the nonintegral vertex in the latter case.
\end{enumerate}

\item A reflexive polygon need not define a smooth toric surface.
For each of the following polygons, compute the determinants of
consecutive primitive ray generators of its normal fan:
\[
\Conv\{(1,0),(0,1),(-1,-1)\},
\qquad
\Conv\{(1,0),(0,1),(-1,-2)\}.
\]
Determine which associated toric surfaces are smooth. Explain why
both polygons can nevertheless be reflexive.

\item Let \(X\) be a smooth toric del Pezzo surface and let
\(\widetilde X\to X\) be the blowup of a torus fixed point.
Show directly from the corner-cut picture that
\[
(-K_{\widetilde X})^2=(-K_X)^2-1.
\]
Then explain why blowing up a fixed point of
\(\Blp{\PP^2}{p,q,r}\) produces a surface whose anticanonical divisor
is not ample.

\item For each of the five smooth toric del Pezzo polygons in
\cref{fig:five-smooth-toric-del-pezzo-polygons}, compute the polar
polygon and verify experimentally that
\[
|\partial P\cap M|+|\partial P^\circ\cap N|=12.
\]
Record how the two summands change under each toric blowup. Do not try
to prove the identity in general yet; formulate instead what a proof
would need to explain.

\end{exercises}

\clearpage
\section*{Challenge exercise}
\addcontentsline{toc}{section}{Challenge exercise}
\markboth{Challenge exercise}{Challenge exercise}

\noindent
\textbf{A crepant resolution hidden in a reflexive triangle.}

\medskip

Let \(N=\ZZ^2\), and consider the complete fan \(\Sigma\) with rays
\[
v_0=(1,0),
\qquad
v_1=(-1,2),
\qquad
v_2=(0,-1),
\]
and maximal cones generated by each consecutive pair. The positive
relation
\[
v_0+v_1+2v_2=0
\]
identifies the corresponding toric surface with
\[
Y=X_\Sigma\simeq\PP(1,1,2).
\]

The cone \(\Cone{v_0,v_1}\) is singular. Insert the primitive ray
\[
w=(0,1)
\]
and let \(\widetilde\Sigma\) be the resulting subdivision. Write
\[
\pi:X_{\widetilde\Sigma}\longrightarrow Y
\]
for the induced proper birational toric morphism.

The goal is to prove that \(\pi\) is a \emph{crepant resolution}:
\[
X_{\widetilde\Sigma}\text{ is smooth},
\qquad
K_{X_{\widetilde\Sigma}}=\pi^*K_Y.
\]
You will also see why resolving the singularity changes the
anticanonical divisor from ample to merely basepoint free.

\begin{figure}[ht]
\centering
\begin{tikzpicture}[scale=1.05]
  % polar triangle
  \begin{scope}[xshift=-3.6cm]
    \draw[->,black!55] (-1.7,0)--(2.05,0)
      node[below right,xshift=2pt] {\(N_\RR\)};
    \draw[->,black!55] (0,-1.55)--(0,2.35);
    \filldraw[fill=DictionaryBlue!12,draw=DictionaryBlue,very thick]
      (1,0)--(-1,2)--(0,-1)--cycle;
    \fill[WarningRed] (0,1) circle (2.7pt);
    \fill (1,0) circle (2.2pt);
    \node[right] at (1,0) {\(v_0\)};
    \fill (-1,2) circle (2.2pt);
    \node[above left] at (-1,2) {\(v_1\)};
    \fill (0,-1) circle (2.2pt);
    \node[below] at (0,-1) {\(v_2\)};
    \node[WarningRed,right] at (0,1) {\(w\)};
    \node[font=\bfseries] at (0,2.75) {the reflexive polar triangle};
    \node[font=\small,align=center] at (0,-2.0)
      {the new ray is a lattice point\\on the edge \([v_0,v_1]\)};
  \end{scope}

  % refined fan
  \begin{scope}[xshift=3.65cm]
    \fill[DictionaryBlue!9] (0,0)--(1.45,0)--(0,1.45)--cycle;
    \fill[SurfaceGreen!10] (0,0)--(0,1.45)--(-1.15,2.3)--cycle;
    \fill[ExampleGold!13] (0,0)--(-1.15,2.3)--(0,-1.4)--cycle;
    \fill[PausePurple!10] (0,0)--(0,-1.4)--(1.45,0)--cycle;
    \draw[->,ray] (0,0)--(1.55,0) node[right] {\(v_0\)};
    \draw[->,ray] (0,0)--(-1.15,2.3) node[left] {\(v_1\)};
    \draw[->,ray] (0,0)--(0,-1.55) node[below right] {\(v_2\)};
    \draw[->,WarningRed,very thick] (0,0)--(0,1.65)
      node[above right] {\(w\)};
    \draw[black!35,dashed] (1.45,0)--(-1.15,2.3);
    \node[font=\bfseries] at (0.2,2.95) {the refined fan};
    \node[font=\small,align=center] at (0.15,-2.35)
      {the singular cone is split\\into two smooth cones};
  \end{scope}
\end{tikzpicture}
\caption{The midpoint \(w\) of the polar edge \([v_0,v_1]\) lies on
the same height-one hyperplane as its endpoints. Adding its ray
resolves the singular cone without changing the anticanonical
polytope.}
\label{fig:crepant-resolution-p112-challenge}
\end{figure}

\begin{enumerate}[label=\textbf{Hint \arabic*.},wide=0pt]

\item \textbf{Locate the singularity.}

Compute the determinants of the three maximal cones of \(\Sigma\).
Show that two cones are smooth and that
\[
|\det(v_0,v_1)|=2.
\]
Conclude that \(Y\) has one cyclic quotient singularity of order two.
Why is the fan nevertheless complete?

\item \textbf{Recover the reflexive anticanonical triangle.}

Write \(m=(a,b)\in M_\RR\). Show that
\[
P_{-K_Y}
=
\{(a,b):a\geq-1,\ b\leq1,\ a-2b\leq1\}.
\]
Find its three vertices and verify that it is a lattice polygon with
unique interior lattice point \(0\). Compute its polar and show that
\[
P_{-K_Y}^\circ
=
\Conv\{v_0,v_1,v_2\}.
\]

\item \textbf{Check the Gorenstein condition locally.}

On the singular cone \(\tau=\Cone{v_0,v_1}\), find an element
\(m_\tau\in M\) satisfying
\[
\langle m_\tau,v_0\rangle
=
\langle m_\tau,v_1\rangle
=1.
\]
Show that one may take
\[
m_\tau=(1,1).
\]
Conclude from the height-one criterion that \(K_Y\) is Cartier near
the singular point. Explain why this is the local reason that
\(Y\) is Gorenstein rather than merely \(\QQ\)-Gorenstein.

\item \textbf{Resolve the cone.}

Observe that
\[
w=\frac{v_0+v_1}{2}=(0,1).
\]
Compute
\[
\det(v_0,w),
\qquad
\det(w,v_1),
\]
and show that the subdivided fan \(\widetilde\Sigma\) is smooth.
Compare its four rays with the standard fan of \(\mathbb F_2\), and
deduce
\[
X_{\widetilde\Sigma}\simeq\mathbb F_2.
\]

\item \textbf{Why the resolution is crepant.}

Check that the new ray lies on the same height-one hyperplane:
\[
\langle m_\tau,w\rangle=1.
\]
Use the Cartier data for \(K_Y\) to compute the coefficient of the
new divisor \(E=D_w\) in \(\pi^*K_Y\). Compare with the toric formula
\[
K_{X_{\widetilde\Sigma}}
=-D_{v_0}-D_{v_1}-D_{v_2}-E
\]
and prove
\[
\boxed{K_{X_{\widetilde\Sigma}}=\pi^*K_Y.}
\]
This equality is the definition of a crepant resolution.

\item \textbf{The anticanonical polygon does not change.}

The new ray contributes the additional inequality
\[
\langle(a,b),w\rangle=b\geq-1.
\]
Show that this inequality is already implied by
\[
a\geq-1,
\qquad
a-2b\leq1.
\]
Conclude that
\[
P_{-K_{X_{\widetilde\Sigma}}}
=
P_{-K_Y}.
\]
Interpret this equality as
\[
-K_{X_{\widetilde\Sigma}}
=
\pi^*(-K_Y).
\]

\item \textbf{Ample downstairs, only nef upstairs.}

Since \(P_{-K_Y}\) has normal fan \(\Sigma\), the divisor \(-K_Y\)
is ample. Explain why the refined fan \(\widetilde\Sigma\) is strictly
finer than the normal fan of the same polygon. Deduce that
\[
-K_{X_{\widetilde\Sigma}}
\]
is basepoint free but not ample.

Equivalently, show that the support function is linear across the new
wall generated by \(w\). Which curve is contracted by the associated
anticanonical morphism?

\item \textbf{See the exceptional curve numerically.}

The neighboring rays of \(w\) satisfy
\[
v_0+v_1=2w.
\]
Use the smooth toric surface intersection formula to show that
\[
E^2=-2.
\]
Then compute
\[
(-K_{X_{\widetilde\Sigma}})\cdot E=0.
\]
Explain why this confirms both that \(-K\) is not ample and that the
anticanonical morphism contracts exactly the \((-2)\)-curve \(E\).

\item \textbf{The general pattern.}

Summarize the example in the following three languages:
\[
\begin{array}{c|c|c}
\textbf{Fan} & \textbf{Divisors} & \textbf{Polytope}\\
\hline
\text{insert a height-one ray}
&
K_{\widetilde X}=\pi^*K_X
&
\text{add a redundant distance-one inequality}
\\
\text{refine one cone}
&
-K_{\widetilde X}\text{ loses ampleness}
&
P_{-K}\text{ stays unchanged}
\\
\end{array}
\]
Why is this phenomenon special to rays lying on the height-one slice?
What discrepancy would you expect if the new ray lay above or below
that slice?

\end{enumerate}

\begin{tcolorbox}[
  colback=PredictionTeal!5,
  colframe=PredictionTeal,
  title={The lesson of the resolution},
  fonttitle=\bfseries
]
A crepant toric resolution does not create a new anticanonical
polytope. It refines the fan by inserting lattice points on the
height-one boundary of the polar polytope. The canonical class is
preserved, but strict convexity may be lost:
\[
\boxed{
\text{Gorenstein Fano downstairs}
\quad\longrightarrow\quad
\text{basepoint-free }-K\text{ upstairs}.
}
\]
For \(\PP(1,1,2)\), the resolution is \(\mathbb F_2\), and the new
ray produces the \((-2)\)-curve contracted by the anticanonical map.
\end{tcolorbox}

\clearpage
\section*{Notes and references}
\addcontentsline{toc}{section}{Notes and references}
\markboth{Notes and references}{Notes and references}

Invariant differential forms and the toric residue sequence are
developed in Cox--Little--Schenck, \S8.1. The canonical divisor and
the Cartier and Gorenstein criteria are treated in \S8.2. Fano toric
varieties, weighted projective examples, and the correspondence with
reflexive polytopes appear in \S8.3.

Our polar convention is
\[
P^\circ
=
\{u:\langle m,u\rangle\geq-1\text{ for every }m\in P\}.
\]
This is adapted to the toric divisor inequalities used throughout the
notes. With the opposite sign convention, some sources write the
polar using \(\leq1\); the geometry is the same after replacing one
polytope by its negative.

The classification of smooth complete toric surfaces is developed in
Cox--Little--Schenck, Chapter~10. Combining that classification with
the anticanonical positivity criterion gives the five smooth toric
del Pezzo surfaces displayed in \cref{thm:classification-smooth-toric-del-pezzo}.
The sixteen reflexive polygons include singular Gorenstein del Pezzo
surfaces as well.

The challenge exercise illustrates a basic principle of toric
birational geometry. Subdividing a Gorenstein cone by rays lying on
its height-one slice is crepant. In the example
\[
\mathbb F_2\longrightarrow\PP(1,1,2),
\]
the exceptional curve has self-intersection \(-2\), the canonical
class pulls back unchanged, and the anticanonical polygon remains the
same. This picture will reappear when resolutions and birational
transformations are studied more systematically.

The boundary identity
\[
|\partial P\cap M|+|\partial P^\circ\cap N|=12
\]
for reflexive polygons is traditionally called the twelve-point
theorem. We have used it only as an experimentally visible prediction.
A later Riemann--Roch discussion will explain why the number twelve is
natural from the geometry of smooth toric surfaces.
